\documentclass[11pt,a4paper]{article}
\usepackage[margin=2.2cm]{geometry}
\usepackage[T1]{fontenc}
\usepackage{booktabs}
\usepackage{array}
\usepackage{enumitem}
\usepackage{xcolor}
\usepackage{listings}
\usepackage{hyperref}
\hypersetup{colorlinks=true,linkcolor=blue!50!black,urlcolor=blue!50!black}

\lstset{
  language=C++,
  basicstyle=\ttfamily\small,
  keywordstyle=\color{blue!60!black}\bfseries,
  commentstyle=\color{green!40!black}\itshape,
  stringstyle=\color{red!50!black},
  columns=fullflexible,
  breaklines=true,
  frame=single,
  rulecolor=\color{black!25},
  backgroundcolor=\color{black!3},
  xleftmargin=4pt,xrightmargin=4pt,
  showstringspaces=false
}
\newcommand{\code}[1]{\texttt{\detokenize{#1}}}
\setlist{itemsep=2pt,topsep=4pt}

\title{Changelog: v20 $\rightarrow$ v25.1.0\\[2pt]\large v21: survival search, free steps, routing, landing safety, rushes and the king's plan;\\ v22: the judge's points budget, the queen, the centre rush; v23: per-map settings and the queen's survival; v24: portals; v25: dispersion; v25.1.0: queen wait splits}
\author{}
\date{}

\begin{document}
\sloppy
\maketitle
\vspace{-2.5em}

\section*{Summary}
\textbf{Latest entry: v25.1.0 (6 October 2026).} Based on \code{bot-v25.0.0-omar}, this version changes only the queen's
emergency split: she detaches two segments, waits in place and reassesses movement next turn. The focused tests passed;
six sandbox matches finished 3--3 with no runtime errors. See section~\ref{sec:queenwait251} for the behavior, limits and evidence.

v21 was built in five rounds.

\textbf{Part I, the survival search.} A depth-first \emph{survival search} (\code{survival.hpp}) asks, for each
candidate move, whether our dragon can keep moving for its own length plus two more turns, playing out the
engine's movement rules exactly. Its result feeds the safety filter (\code{critical_ok}) and the position score.
Two supporting changes came out of testing: traps are priced by whether a split-escape is actually possible,
and workers of length $\geq 8$ may use the reserved escape slot at the unit cap.

\textbf{Part II, the free-steps rule and routing.}
\begin{itemize}
      \item The game update's sprint rule ($\lceil L/4\rceil$ free steps a turn) is now modelled everywhere the bot
            reasons about movement: its own moves, enemy reach and the simulator.
      \item The bot uses the rule itself: besides single steps it weighs \emph{free moves} of several steps.
      \item Pearl routing (\code{route.hpp}) is rebuilt around four ideas:
            \begin{itemize}
                  \item time-aware paths: bodies are passable once they have moved on;
                  \item arrival measured in turns rather than steps;
                  \item who gets to each pearl first, by real paths from every visible head;
                  \item a survival search from where each route ends, so that a route that would box us in, above
                        all the king, is never chosen.
            \end{itemize}
      \item Testing showed free moves could carry the king to the edge of its vision, next to enemies it could not see,
            so valuable dragons now price that exposure.
\end{itemize}
\textbf{Part III, landing safety and teammate memory.} Free moves never leave the 7$\times$7 window, but they can
\emph{end} at its edge. Long snakes then got boxed in, mostly by allies they could not see, 5--10$\times$ more often than
after a single step. The changes:
\begin{itemize}
      \item bodies are remembered for as long as they will really be there;
      \item allies' moves are modelled from what we see and from their status packets;
      \item long snakes pay a price for landing next to tiles they cannot see;
      \item the king no longer outruns its feeders in the late game.
\end{itemize}

\textbf{Part IV, going back for remembered pearls.} A long king that left pearls behind, such as pearls 30 steps
away around its own body, never went back: routes only looked 20 steps ahead, which is one turn for a king of 80.
The changes:
\begin{itemize}
      \item routes now reach a few turns of free steps;
      \item a \emph{rush} follows the route for all our free steps, also over ground we cannot see now but have seen,
            priced by how likely it is that something has moved there since;
      \item bodies running out of sight are predicted where they continue.
\end{itemize}
This round also found and fixed a bug that predates v21. Dragon ids pass 4096 in long games and our 12-bit status
packets wrapped them, so feeders lost track of the king and teams were wiped out.

\textbf{Part V, the king's own moves.} Three changes for the king and growers:
\begin{itemize}
      \item the survival search now proves a way on for their whole length, not 20 moves;
      \item a dragon that does not know all of its body treats unseen tiles next to where it ends as possibly its own body;
      \item each of the king's best moves is weighed by a search over its own multi-step moves for the next two turns.
\end{itemize}
Measuring showed that kings rarely die to their own body alone. They get boxed in by other dragons and then
split-escape about once in ten turns, often in chains (see \S\ref{sec:cascade}).

The queens/win-condition update is \textbf{not} implemented.

\textbf{v22 (Part VI), the judge's points budget.} The judge meters every dragon's turn in CPU points (100M a turn); a
turn over it gets no reply, and the dragon dies with ``no valid action''. Measured with the judge's own pricing, v21's
long dragons went over on hundreds of turns a game (kings averaged 164M a turn, up to 1.2 billion). v22 reads its own
spending off the judge's clock, which counts points, and stops each search at its phase's deadline. The same searches
were also made cheaper: a survival-search position by 2.5$\times$, the king's plan options by over 10$\times$. No measured
turn of v22 goes over 87M, and v22 keeps v21's lookahead wherever it fits: the budget changes the chosen move on 0\% of
turns below length 100 and about 8\% at 150+.

\begin{center}
      \begin{tabular}{@{}>{\raggedright}p{4.2cm}p{11.4cm}@{}}
            \toprule
            File                    & Change                                                                                                                                                                                                                                                                                                                                                                                                                                                                                                                        \\
            \midrule
            \code{survival.hpp}     & \textbf{New}: the survival DFS (Part I); bodies-only timings for routing, free-step head reach, \code{survival_from} for route ends (Part II); remembered bodies' lifetimes, allies' known lengths and status packets (Part III); deep searches, unknown-body caution, the king's plan (Part V)                                                                                                                                                                                                                               \\
            \code{route.hpp}        & \textbf{Rewritten} (Part II): time-aware searches, ETAs from every visible head, turn-based route values, BFS-guided replay with a route-end survival check, free moves; no exploring free moves for the late-game king (Part III); route range in turns, rushes (Part IV); deep route-end checks for the king (Part V)                                                                                                                                                                                                       \\
            \code{main.cpp}         & Movement-rule helpers; config; \code{Candidate} fields; turn order; safety filters; escape slot; \code{simulate_path}, threat maps, reach threats, sprints and the ally contest under the new rule; enemy-density tracking; body memory (\code{busy_until}, \code{remember_bodies}), \code{known_length}, ally density (Part III); fill rate and occupancy, remembered tiles in \code{simulate_path}, hidden segments, 12-bit ids (Part IV); deep-search and plan constants, \code{Candidate::plan_value}, plan hook (Part V) \\
            \code{scoring.hpp}      & \code{trap_cost}; new \code{position_value} branches; goal progress for moves of $k$ steps; route progress; fight reach; \code{unseen_risk} (Part III: a box-in part for every dragon of 10+)                                                                                                                                                                                                                                                                                                                                 \\
            \code{explore.hpp}      & Pearl sprints: survival search, new rule, contest by turns; \code{disperse_score} per step                                                                                                                                                                                                                                                                                                                                                                                                                                    \\
            \code{endgame.hpp}      & Death-yield comment; ``cheap way out'' by paid steps; \code{feeding_score} per step; 12-bit ids (Part IV)                                                                                                                                                                                                                                                                                                                                                                                                                     \\
            \code{coordination.hpp} & 12-bit ids (Part IV)                                                                                                                                                                                                                                                                                                                                                                                                                                                                                                          \\
            \code{helper.hpp}       & Replaced with the upgraded-protocol release (still \code{PROTOCOL_MAJOR = 3}). Against ours, only the \code{make_moves} comment differs: it now states the free-steps rule (the first $\lceil L/4\rceil$ steps free, at least 2 segments kept), which the bot and \code{sim/engine.py} already follow. The compiled bot is byte-identical.                                                                                                                                                                                    \\
            \code{sim/engine.py}    & The free-steps rule (\code{FREE_MOVES} switch)                                                                                                                                                                                                                                                                                                                                                                                                                                                                                \\
            \code{sim/} probes      & \code{probe_aggr.py} (metrics), new \code{probe_trapped.py}, \code{probe_bigdeath.py}, \code{probe_turn.py}, \code{scen_route.py}, \code{probe_freemove.py}, \code{probe_hitbody.py}, \code{ab_trap.py}, \code{probe_fill.py}; \code{scen_route.py} gains \code{loopback} and \code{ownpocket}; \code{probe_selfdeath.py}, \code{probe_kingtrap.py}, \code{probe_ksplit.py}, \code{probe_kroom.py}, \code{probe_kingself.py}, \code{probe_bodyk.py}                                                                           \\
            \bottomrule
      \end{tabular}
\end{center}

\tableofcontents

%=====================================================================
\part{The survival search}
%=====================================================================
\section{Movement rules the search is built on}
These were confirmed before writing the search and are what \code{surv_dfs} simulates:
\begin{enumerate}
      \item \textbf{Self-collision includes the tail}, and is checked \emph{before} the tail advances.
            Within a multi-step move, each step is checked against the body as the previous step left it.
      \item \textbf{Dragons act one at a time}, and each sees the actual board at its turn.
      \item \textbf{Sprints} check collisions at every step and remove one extra tail segment per paid step
            (Part II: the first $\lceil L/4\rceil$ steps of a turn are free).
      \item \textbf{Pay before eat}: the validity check for a step happens before that step's pearl is eaten,
            so a pearl does not save a dragon on its death step.
      \item \textbf{Death yield}: a dragon of length $L$ always drops $\lceil L/2 \rceil$ pearls.
      \item \textbf{No simultaneous entry}: two dragons never move into a tile at the same time.
\end{enumerate}

\section{New file: \texttt{survival.hpp}}
Included inside \code{ReplayDerivedBot}, before \code{scoring.hpp}.

\subsection{What it computes}
\begin{lstlisting}
struct SurvivalResult {
    int depth = -1;        // most moves survived on any path (-1: not searched)
    int target = 0;
    bool complete = false; // every path searched (budget did not cut it short)
    bool ok()   const { return depth >= 0 && depth >= target; }
    bool dead() const { return depth >= 0 && complete && depth * 4 < target * 3; }
};
\end{lstlisting}
\begin{itemize}
      \item \textbf{Target depth}: $\min(L + \code{DFS_EXTRA_DEPTH},\ \code{DFS_MAX_DEPTH}) = \min(L+2, 20)$,
            where $L$ is our length after the candidate move. Surviving $L$ moves means our whole
            current body has moved off its tiles; the $+2$ is margin.
      \item \textbf{ok}: some path reached the target.
      \item \textbf{dead}: the search finished in full and the best path is short of $3/4$ of the target. This is a trap.
      \item Anything else (including a budget cut) is ``unknown'', and the old v20 heuristics decide.
      \item Our own moves are searched \emph{one step a turn}. Free steps are optional, so we can always go that slowly,
            and that pace gives other bodies the most turns to move on.
\end{itemize}

\subsection{State (new members)}
\begin{itemize}
      \item \code{other_free_at_}: per tile, how many moves the other dragons must have made before the tile is free (\code{INF}: never).
      \item \code{body_free_at_} (Part II): the same for \emph{bodies only}, meaning other dragons' segments and remembered bodies.
            Our own segments and the head-reach margin are left out. Routing and the ETA searches use it.
      \item \code{surv_occ_}: our body during the search (occupancy counts).
      \item \code{surv_eaten_}: pearls eaten along the current path, so the same pearl is not eaten twice.
      \item \code{surv_nodes_}, \code{surv_turn_nodes_}: per-candidate and per-turn node counters.
      \item Part II: \code{surv_route_nodes_} (the route checks' own per-turn counter), plus \code{surv_counter_}, \code{surv_cap_}
            and \code{surv_offset_}, which select the counter, cap and time offset for the current search.
\end{itemize}

\subsection{\texttt{prepare\_survival()}: once per turn}
Called from \code{play_turn}. Part II moved it before the objective map and routing (see \S\ref{sec:order}).
It resets both node counters and fills \code{other_free_at_} and \code{body_free_at_}:
\begin{enumerate}
      \item A tile with a \textbf{remembered} (out-of-sight) body blocks for \code{DFS_REMEMBERED_STEPS} $=3$ moves (both maps).
      \item \textbf{Our own unplaceable segments} (\code{own_idx_ == -2}, never seen since birth) free after
            $\max(1,\ L - \code{known_body_})$ moves, since they are in the unknown tail (\code{other_free_at_} only).
      \item \textbf{Every visible part of another dragon} is a wall (\code{INF}) by default (both maps).
      \item For each \textbf{visible head}, the body is traced (\code{trace_body}). Segment $i$ frees after
            $\max(\code{DFS_OTHER_MIN_WAIT},\ \mathit{est} - i + 1)$ moves, where $\mathit{est}$ is that
            dragon's estimated length (both maps). Each dragon moves once between two of ours; the $+1$ covers it eating
            (its tail then stays); and nothing frees before our second move. Segments whose head we cannot see stay walls.
      \item \textbf{Head reach}: tiles a visible head could be on before our next move get a value of at least 2, so they are
            not counted on for our first step. In Part I these were its non-reverse next tiles. In Part II they are every tile
            within \code{free_moves(est)} steps (\code{tiles_within}), since its body follows it there. This applies to
            \code{other_free_at_} only; where heads go after that is left to the fight search.
\end{enumerate}

\subsection{\texttt{survival(pm, heading)}: per candidate}
\begin{itemize}
      \item Builds our body after the move, head first: \code{pm.path} reversed, then \code{own_trail_},
            truncated to \code{pm.length}. Any remaining length is an \code{unknown} count of tail segments we cannot place.
      \item Part II: the search itself is now \code{survival_from(body, unknown, heading, length, offset, route)}, which
            \code{survival} calls with offset 0 and the moves' budget.
      \item If the per-turn budget is already spent, it returns ``not searched'' (\code{depth = -1}).
\end{itemize}

\subsection{\texttt{surv\_dfs}: the search}
\begin{itemize}
      \item Stops with success at the target. It stops with \code{cut = true} when either budget runs out:
            \code{DFS_NODE_BUDGET} $=4000$ per search, and \code{DFS_TURN_BUDGET} $=40000$ per turn for moves.
            Route checks have their own \code{ROUTE_DFS_TURN_BUDGET} $=24000$.
      \item \textbf{Children} are the (up to three) non-reverse directions whose destination is passable and known.
            The tile must also hold none of our body (tail included, per rule 1) and have
            $\code{other_free_at_} \leq \mathit{step} + \code{surv_offset_}$. Kelp blocks, the map wraps, and known portals are
            followed. Unknown portals are not taken, since we cannot see where they lead.
      \item \textbf{Ordering}: children are sorted by \emph{openness}, the count of free neighbours at the next step,
            most open first. In open water the first path reaches the target, so the search is cheap; the budget only bites in tight spots.
            This is a manual insertion sort over at most three elements (see \S\ref{sec:iter}).
      \item \textbf{Moving}: the head is pushed. If the tile has a pearl seen this round that is not already eaten on this path,
            we grow and the tail stays. Otherwise an unknown tail segment is consumed if any remain, else the real tail tile is popped.
      \item \textbf{Backtracking} undoes every change: the head push, the tail pop, the unknown count and the eaten marks.
      \item Records the deepest depth reached in \code{best}.
\end{itemize}

\section{Part I changes outside \texttt{survival.hpp}}
\subsection{Configuration (\texttt{main.cpp})}
\begin{lstlisting}
#ifndef LONG_WORKER_ESCAPE
#define LONG_WORKER_ESCAPE 8
#endif
constexpr int LONG_WORKER_ESCAPE_LENGTH = LONG_WORKER_ESCAPE;

// v21: survival search
constexpr int DFS_EXTRA_DEPTH = 2;      // target = length + this ...
constexpr int DFS_MAX_DEPTH = 20;       // ... capped here
constexpr int DFS_NODE_BUDGET = 4000;   // per search
constexpr int DFS_TURN_BUDGET = 40000;  // per turn (moves)
constexpr int DFS_REMEMBERED_STEPS = 3; // remembered body blocks this long
constexpr int DFS_OTHER_MIN_WAIT = 2;   // no other segment frees before this
constexpr double P_DFS_SHORT = 0.9;     // trap probability per unit shortfall
\end{lstlisting}
\code{LONG_WORKER_ESCAPE} sits after \code{ESCAPE_UNIT_RESERVE = 2}. Its comment records why it exists:
boxed-in workers at the unit cap used to die whole, about 140 deaths of dragons of length 10+ per 200 rounds on a big map.
Compiling with \code{-DLONG_WORKER_ESCAPE=1000} turns the feature off.

\subsection{\texttt{Candidate} fields}
\begin{lstlisting}
int dfs_depth = -1, dfs_target = 0; bool dfs_complete = false;
bool dfs_ok() const   { return dfs_depth >= 0 && dfs_depth >= dfs_target; }
bool dfs_dead() const { return dfs_depth >= 0 && dfs_complete
                               && dfs_depth * 4 < dfs_target * 3; }
\end{lstlisting}

\subsection{Include}
\code{#include "survival.hpp"} is added before \code{scoring.hpp}.

\subsection{Escape slot policy (\texttt{escape\_slot\_available})}
Previously a worker always had to leave \code{ESCAPE_UNIT_RESERVE} $=2$ unit slots free before it could
split-escape. Now:
\begin{lstlisting}
reserve = king ? 0
        : (valuable_now_ || length >= LONG_WORKER_ESCAPE_LENGTH) ? 1
        : ESCAPE_UNIT_RESERVE;
\end{lstlisting}
A worker of length $\geq 8$ may use the second-to-last slot to split out of a trap, like a valuable dragon.
Losing 8+ segments whole costs far more than one extra unit near the cap.

\subsection{Safety filters}
In \code{evaluate_single}, after the room facts are gathered, \code{survival(pm, direction)} runs and is stored in the
candidate. The room part of \code{critical_ok} is now:
\begin{lstlisting}
room_ok = dfs_ok()
       || (!dfs_dead() && exits > 0
           && space >= pm.length + unknown_body() + 2 && !trap);
\end{lstlisting}
\begin{itemize}
      \item A move the search proves survivable is safe even when the flood-fill heuristics call it tight.
      \item A move the search proves a trap is unsafe even when they look fine.
      \item Unknown falls back to v20.
      \item The other conditions are unchanged: known destination, no unknown portal, no head trade, no occupied exit,
            and for valuable dragons no hard threats or blind portal.
\end{itemize}
The same rule, run from each sprint's end, decides \code{critical_ok} for escape sprints
(with \code{hard_threats == 0}) and pearl sprints (with \code{!valuable || hard == 0}).

\subsection{Scoring (\texttt{scoring.hpp})}
\textbf{\code{trap_cost}.} In v20 a trap was priced as if a split-escape were always available (2 segments plus a unit). At the unit
cap it is not, and the dragon dies whole, so traps were underpriced by about $3\times$. Now:
\begin{lstlisting}
bool const can_escape = L >= 4 && escape_slot_available()
                        && controller_.can_split(L - 2);
return can_escape ? seg(2.0 + Config::UNIT_VALUE_OURS) : death_cost();
\end{lstlisting}
\textbf{\code{position_value}}: new branches after the pocket-plan case.
\begin{itemize}
      \item \code{dfs_ok()}: the room is proven. Only the \code{P_ONE_EXIT} risk is charged when \code{exits == 1};
            the space and dead-end risks are skipped.
      \item Searched in full but short: \code{v -= risk(P_DFS_SHORT * shortfall, trap_cost())}, where
            $\mathit{shortfall} = (\mathit{target} - \mathit{depth})/\mathit{target}$.
      \item Not searched, or cut short by the budget: the v20 exits/space/dead-end risks, unchanged.
\end{itemize}

\subsection{Debug log and comments}
\begin{itemize}
      \item The \code{CAND} debug line ends with \code{dfs <depth> <target> C|-}, where \code{C} means the search was complete.
      \item \code{endgame.hpp}: ``Full yield'' was misleading, so the comment now reads:
            \emph{Every death leaves $\lceil L/2\rceil$ pearls on the body's tiles; dying on our own neck here (rather than in a trade)
                  puts them where the king will eat them.}
\end{itemize}

\section{How the search got here}\label{sec:iter}
\begin{enumerate}
      \item \textbf{First version (optimistic).} Other dragons' tails freed at step 1 and other heads were ignored.
            Since \code{dfs_ok()} drops the \code{exits == 0} risk, dragons stepped onto tiles with no free neighbour.
      \item \textbf{Fix: cautious timing.} Added \code{DFS_OTHER_MIN_WAIT = 2} and the $\mathit{est} - i + 1$ eating margin,
            and blocked other heads' next tiles for our first step.
      \item \textbf{Fix: trap pricing.} Nearly all self-deaths were \code{TRAPPED_SELF} (no legal move)
            at $\geq 62$ units, where workers cannot split-escape. \code{trap_cost} now checks \code{escape_slot_available() && can_split}.
      \item \textbf{Escape slot.} Even correctly priced traps still happen at the cap, so long workers were allowed the reserve slot.
      \item \textbf{Compiler warning.} \code{std::sort} on \code{std::array<Child,3>} triggered a GCC
            \code{-Warray-bounds} false positive. It was replaced with a manual insertion sort.
\end{enumerate}

\section{Part I results (old movement rule)}
No-legal-move deaths in 200 rounds on seed 20 against the baseline:
\begin{center}
      \begin{tabular}{@{}lrr@{}}
            \toprule
            Version             & Trapped deaths & of which length 10+ \\
            \midrule
            v20                 & 318            & 157                 \\
            v21, cautious DFS   & 297            & --                  \\
            v21, + trap pricing & 270            & --                  \\
            v21, + escape slot  & 238            & 92                  \\
            \bottomrule
      \end{tabular}
\end{center}
\begin{itemize}
      \item Walls around trapped heads were mostly allies (ally body 463--499, ally head 86--111), then own body ($\sim$450),
            kelp ($\sim$120) and enemies ($\sim$25).
      \item Full games against v20: 2--2 without the escape slot; \textbf{4--2} with it, with accidental self-deaths about
            27\% lower on big maps.
      \item The scenario tests passed and the coil stress test failed exactly as v20 did; the maximum turn time was 5--13\,ms.
\end{itemize}

%=====================================================================
\part{The free-steps rule, free moves and routing}
%=====================================================================
\section{The free-steps rule}\label{sec:rule}
\subsection{The rule as modelled}
In one turn a dragon of length $L$ may take $\lceil L/4\rceil$ steps for free. Every step after those costs a segment
(its tail loses one more) and is allowed only while the dragon is longer than 2. That gives 1 free step at lengths 1--4
(the old rule), 2 at 5--8, 3 at 9--12 and 5 at 20.

\textbf{Assumption}: the allowance is per turn, and $L$ is the length at the \emph{start} of the turn.

\begin{lstlisting}
constexpr int free_moves(int length) { return std::max(1, (length + 3) / 4); }
constexpr int paid_steps(int length, int steps) { return std::max(0, steps - free_moves(length)); }
// The most steps it can take in one turn (eating nothing on the way).
constexpr int max_turn_steps(int length) { return free_moves(length) + std::max(0, length - 2); }
// Turns to take `steps` steps at full free speed.
constexpr int turns_for(int steps, int length);   // ceil(steps / free_moves(length))
\end{lstlisting}

\subsection{Where the bot now uses it}
\begin{center}
      \begin{tabular}{@{}>{\raggedright}p{4.3cm}p{11.3cm}@{}}
            \toprule
            Place                     & Change                                                                                                                    \\
            \midrule
            \code{simulate_path}      & The pay check applies only to paid steps (\code{s >= free && length_now <= 2}). Length before step $s$ is
            $L + \mathit{eaten} - \max(0, s - \mathit{free})$; after the move, $L + \mathit{eaten} - \code{paid_steps}$.                                          \\
            \code{play_turn}          & The projected length sent in the status packet subtracts \code{paid_steps}.                                               \\
            Reach threats (valuable)  & Hard reach is \code{max_turn_steps(lower_length)} (was \code{lower_length - 1}). Soft reach for a truncated body is
            \code{max_turn_steps} of the larger of the typical enemy length and that dragon's own estimate (what we see $+2$), still capped by
            \code{HIDDEN_BODY_REACH}.                                                                                                                             \\
            \code{build_threat_maps}  & An enemy head threatens every tile within its free steps (\code{tiles_within}); an ally's next step is still what
            crowds us. For a reach of 1 this is the same set as before.                                                                                           \\
            Fight search              & An enemy's sprint into our head reaches \code{clamp(max_turn_steps(len), 1, max(FIGHT_SPRINT_REACH, free_moves(len)))}.   \\
            Survival search           & Head reach covers its free steps (Part I, \code{prepare_survival}).                                                       \\
            Tactical (attack) sprints & Up to \code{min(max_turn_steps(L), MAX_ATTACK_STEPS)} steps; the tie-break cost is per \emph{paid} step.                  \\
            Escape sprints            & Up to \code{min(max_turn_steps(L) - 1, MAX_ESCAPE_STEPS)} steps; the score is \code{(eaten - paid) * SEG}.                \\
            Pearl sprints             & Only for pearls beyond our free steps (closer ones are free moves). Net is eaten minus paid. ``Contested'' means an enemy
            needs no more turns than we would at free speed (by ETA, \S\ref{sec:eta}).                                                                            \\
            Split escape              & A ``cheap way out'' is any safe move paying at most 2 segments (was: at most 3 steps).                                    \\
            Simulator                 & \code{engine.py}: \code{FREE_MOVES = True} applies the rule (pay check and extra tail segment only past the free steps).  \\
            \bottomrule
      \end{tabular}
\end{center}

\section{Free moves}
\subsection{Generation (\texttt{add\_free\_moves}, \texttt{route.hpp})}
Each turn, after the single steps and the tactical and pearl sprints (and before escape sprints), we weigh moves of 2 up to
$\min(\code{free_moves}(L), \code{FREE_MOVE_MAX_STEPS} = 6)$ steps:
\begin{itemize}
      \item A depth-first walk over step sequences (never reversing). Every step is through a tile we can see now, with no portal.
            Each prefix is checked with \code{simulate_path}, which uses the exact engine rules.
      \item One option is kept per (end tile, last direction), the best by a quick look (\code{quick_value}): food on the way,
            exits at the end, and goal progress (along the route by the pearls' own maps).
      \item The best \code{FREE_MOVE_EVALS} $=6$ are evaluated in full by \code{evaluate_path}.
\end{itemize}

\subsection{Evaluation (\texttt{evaluate\_path})}
It gathers the same facts as \code{evaluate_single}: threats at the end, exits, time-aware space and region, known space
(valuable), and the survival search. The move is priced as in \code{move_value}:
\begin{itemize}
      \item goal progress for $k$ steps (half strength in a fight);
      \item every pearl eaten on the way (\code{path_food}), minus paid steps (none);
      \item \code{position_value} at the end and, in a fight, \code{fight_value} from there;
      \item visits to the end tile, teammate spacing (unless it eats), and \code{BLOCKED_SPAWN} for every path tile about to spawn;
      \item the enemy-distance pull.
\end{itemize}
\code{critical_ok} and \code{safe} follow the single-step rules. The new \code{Candidate::free_move} flag makes the
chosen move log as \code{FREE_MOVE}.

\subsection{Goal progress for $k$ steps (\texttt{scoring.hpp})}
\code{goal_progress} now takes the \code{PostMove} instead of a \code{Step}, so the same code prices one step or several.
With $k = 1$ every term is as before.
\begin{center}
      \begin{tabular}{@{}lp{11.6cm}@{}}
            \toprule
            Goal    & Progress                                                                                                                                                        \\
            \midrule
            EAT     & Steps still to go along the route, before minus after, clamped to $\pm k$ (\S\ref{sec:routeprog}); \code{ROUTE_TIE} if the move follows the checked route path. \\
            HUNT    & Grid distance gained to the target (unclamped, as before).                                                                                                      \\
            FEED    & \code{feeding_score(dest, k)}: the gain clamped to $\pm(k+1)$ (was $\pm 2$).                                                                                    \\
            EXPLORE & Objective-map gain capped at $k + 1$ (was 2); \code{disperse_score(dest, k)} clamped to $\pm k$; rendezvous unchanged.                                          \\
            \bottomrule
      \end{tabular}
\end{center}

\section{Routing (\texttt{route.hpp})}
\subsection{Time-aware paths}
\code{timed_bfs(pm, limit, target, out)} gives the fewest steps for our head to reach a tile from where a move \code{pm} leaves us
(or from where we are now). Every body is placed where it will be by then:
\begin{itemize}
      \item Our own segment $j$ (after the move) is passed once the tail has moved past it: at least $\mathit{length} - j + 1$ more steps.
      \item Another dragon's segment is passed once that dragon has made \code{body_free_at_} moves, one a turn. Our step is assumed
            to come in turn $\code{turns_for}(\mathit{step})$, at full free speed, plus the current turn if a move is already
            being made. Going that fast is pessimistic about other bodies.
      \item Kelp, tiles we know nothing of, and marked traps block.
\end{itemize}
It is used in three places:
\begin{itemize}
      \item \code{route_head_dist_}: the map from our head, out to 30 steps (\code{prepare_etas}).
      \item The pearls' own maps (legs between pearls): a body tile is passable if our time-aware map reaches it, since it will
            have moved on by then. The first leg of every route uses \code{route_head_dist_} directly.
      \item Route progress for candidate moves (\S\ref{sec:routeprog}).
\end{itemize}

\subsection{Arrival in turns}
A pearl reached at step $t$ is reached in $\code{turns_for}(t, L)$ turns. The discount (0.92 for keepers, 0.85 for workers)
is now \emph{per turn}; for $L \leq 4$ this is the old per-step rule. Orders that arrive in the same turns are told
apart by a $10^{-4}$-per-step tie-break. Candidate pearls are filtered by time-aware steps within the route radius
(20 for keepers, 14 for workers).

\subsection{Who gets there first}\label{sec:eta}
\code{prepare_etas()} runs once per turn, before the objective map. It computes our time-aware map and one search per
visible head (\code{head_eta_bfs}, 30 steps):
\begin{itemize}
      \item That head's own segments pass once it has moved on.
      \item Our segments pass once we have made that many moves (we move before each of its turns).
      \item Others pass per \code{body_free_at_}.
\end{itemize}
Enemies are merged into \code{enemy_eta_turn_} and \code{enemy_eta_steps_}; allies are kept per head (\code{ally_etas_}).

The results are used in four places:
\begin{itemize}
      \item \textbf{Route values}: a pearl an enemy reaches in fewer turns than our arrival on that route is worth
            \code{ROUTE_CONTEST_LOST} $=0.4$ of itself. Before, any pearl closer to an enemy by \emph{grid} distance than to us by path
            was halved. We win ties: every other dragon's $t$-th turn comes after ours.
      \item \textbf{Keepers}: the fight-risk factor (0.4) now needs an enemy within \code{ROUTE_ENEMY_NEAR} $=4$ \emph{path} steps of the pearl.
      \item \textbf{\code{pearl_taken_by_ally}}: turns at each dragon's own free speed, by time-aware paths for us and visible allies.
            Allies out of sight still use grid distance, now in turns. Ties still go to the lower id, a rule both dragons follow.
      \item \textbf{Pearl sprints}: contested by turns (see \S\ref{sec:rule}).
\end{itemize}

\subsection{Replay and the route-end check}
\code{replay_route} walks a route as the engine would move us, with our body trailing and growing:
\begin{itemize}
      \item Each leg follows a shortest time-aware path from wherever the replay has got to (\code{replay_leg}): a BFS that
            knows the replayed body's segment indices, with the most open first step preferred.
      \item If the growing body gets in the way, the leg is searched again, at most 3 times per leg.
      \item A replay may take at most \code{ROUTE_REPLAY_SLACK} $=8$ steps more than the horizon.
\end{itemize}
From where the route ends, the survival search runs (\code{survival_from}) on the replayed body:
\begin{itemize}
      \item The time offset is $\code{turns_for}(t) - 1$ other-dragon moves; the pearls the route ate are marked eaten.
      \item It uses its own per-turn budget, so routing never starves the move evaluation.
\end{itemize}
The route is accepted only if that search succeeds:
\begin{itemize}
      \item \textbf{King and growers}: \code{ok()}, or cut short by the budget after reaching at least $3/4$ of the target.
      \item \textbf{Workers}: not \code{dead()}.
      \item \textbf{Budget already spent}: the old check (at least one free neighbour) decides.
\end{itemize}
The fallback (the nearest pearl on its own) goes through the same check. The replayed tiles are kept in
\code{route_path_}, its first tile in \code{route_next_cell_}, and the order in \code{route_seq_}.

\subsection{Route progress}\label{sec:routeprog}
\code{route_remaining(pm, limit)} counts the steps still to go after a move:
\begin{itemize}
      \item time-aware steps (\code{timed_bfs}) from where the move leaves us to the first route pearl it has not eaten;
      \item plus the legs between the rest, skipping pearls eaten on the way.
\end{itemize}
\code{route_remaining_now_} is the same from here. EAT progress is the difference, clamped to $\pm k$; a move that eats
the first pearl and carries on toward the second is credited for both. \code{follows_route(pm)} gives \code{ROUTE_TIE}
when the move's tiles are a prefix of \code{route_path_}. The old \code{route_dist()} is removed.

\subsection{Turn order}\label{sec:order}
\begin{lstlisting}
update_explore_target();
prepare_survival(); // v21: when other dragons' bodies move on (survival.hpp)
prepare_etas();     // v21: who reaches which tile first (route.hpp)
build_objective_map();
plan_route();
choose_goal();
...
add_tactical_sprints(candidates);
add_pearl_sprints(candidates);
add_free_moves(candidates); // v21
add_escape_sprints(candidates);
\end{lstlisting}
The debug \code{ROUTE} line now shows \code{rest <route_remaining_now_>} and \code{end_dfs <depth>}.

\section{Valuable dragons: what they cannot see}
A test game showed the danger. A king of length 111 made a 6-step free move to the edge of its window and landed next to
an enemy of length 6 it could not see; the enemy rammed it with a free 2-step move. A single step never ends closer than 2
tiles to the window's edge; a free move can end on it.

\textbf{\code{unseen_risk(pm)}} (\code{position_value}, valuable dragons only) prices this.
\begin{itemize}
      \item Count the tiles within \code{UNSEEN_RADIUS} $=3$ steps of where we end that we cannot see now.
            Every single step has exactly one; an end at the window's edge has up to a dozen.
      \item Each tile is a chance of $\code{P_UNSEEN_ATTACK} = 0.5$ times the enemy heads we see per tile.
            That density is a running average (\code{enemy_density_x1e4_}, updated in \code{observe_heads}).
      \item The total is capped at \code{P_UNSEEN_CAP} $=0.3$, times \code{death_cost()}.
      \item A blind portal exit is not counted here, since portals are priced separately.
\end{itemize}
A second king death (length 85) came from a body running out of sight. It was assumed only as long as a typical enemy,
reaching 3 steps, but it was 8 long and paid 3 segments to reach 5. The soft reach now also uses that dragon's own
estimated length (see \S\ref{sec:rule}).

\section{New and changed constants (Part II)}
\begin{lstlisting}
constexpr int FREE_MOVE_MAX_STEPS = 6, FREE_MOVE_EVALS = 6;
constexpr int UNSEEN_RADIUS = 3;
constexpr double P_UNSEEN_ATTACK = 0.5, P_UNSEEN_CAP = 0.3;
constexpr double ROUTE_CONTEST_LOST = 0.4;
constexpr int ETA_RADIUS = 30;
constexpr int ROUTE_DFS_TURN_BUDGET = 24000;
constexpr int ROUTE_REPLAY_SLACK = 8;
\end{lstlisting}
\code{main.cpp} also gains \code{#include <tuple>}, the four rule helpers, \code{tiles_within} and
\code{enemy_density_x1e4_}. Its header comment adds: \emph{v21 (continued): the ceil(L/4) free-steps rule throughout
      (free\_moves); moves of several free steps; routing with time-aware paths, arrival in turns, who gets there first, and a survival
      check where each route ends (route.hpp).}

\section{Part II testing (new movement rule, against v21 Part I)}
\subsection{Routing scenarios (\texttt{sim/scen\_route.py})}
\textbf{\code{room}.} A king of length 15 next to a fully visible 3$\times$3 walled room with a one-tile door and 3 pearls inside
(eating them leaves 18 segments in 9 tiles), plus 2 pearls in the open.
\begin{itemize}
      \item Both versions eat the open pearls first.
      \item The Part I version then plans a route into the room and spends 4+ turns walking at it; its move-level safety stops it at the door.
      \item The new version rejects every route into the room and goes exploring elsewhere.
\end{itemize}
\textbf{\code{contest}.} A worker of length 6 (2 free steps) and an enemy head 2 tiles from pearl P1 by grid distance, but behind a
kelp wall (30+ steps by path); pearl P2 is 1 step further for us and uncontested.
\begin{itemize}
      \item The Part I version gives P1 up to the enemy and eats nothing in 4 turns.
      \item The new version takes P1 and then heads for P2.
\end{itemize}

\subsection{Games against v21 Part I (both under the new rule)}
\begin{center}
      \begin{tabular}{@{}lllrrl@{}}
            \toprule
            Seed & Map          & New as & New max & Old max & Result                              \\
            \midrule
            21   & 20$\times$36 & A      & 177     & 90      & new wins                            \\
            21   & 20$\times$36 & B      & 174     & 98      & new wins                            \\
            22   & 18$\times$25 & A      & --      & --      & new wins (old eliminated, round 86) \\
            22   & 18$\times$25 & B      & --      & --      & new wins (old eliminated, round 76) \\
            20   & 56$\times$53 & A      & 296     & 97      & new wins                            \\
            20   & 56$\times$53 & B      & 293     & 142     & new wins                            \\
            30   & 44$\times$61 & A      & 220     & 89      & new wins                            \\
            30   & 44$\times$61 & B      & 259     & 100     & new wins                            \\
            \bottomrule
      \end{tabular}
\end{center}
\textbf{8--0.} These games were played before the one-line route-end offset fix (\S\ref{sec:iter2}). Seed 30 re-run with the final
build: new wins both again (max 290 vs 95 as A, 310 vs 136 as B). By round 500 the new build is down to 1--2 units (its workers
have fed the king), against the old build's 9--11. On total length the old build was ahead in both games (638 vs 562, 836 vs 310), so
under the current tiebreak (longest dragon first) the new build wins, but the totals are worth watching once the queen rule
is implemented.

\subsection{Metrics (seed 20, 56$\times$53, 200 rounds)}
\begin{center}
      \begin{tabular}{@{}lrr@{}}
            \toprule
                                                           & New        & v21 Part I \\
            \midrule
            Pearls eaten                                   & 8791       & 4611       \\
            Total length at round 200                      & 1052       & 627        \\
            Segments lost to enemy trades                  & 2382       & 1044       \\
            Our dragons of 16+ hit by enemy-started trades & 46         & 9          \\
            Accidental self-deaths (segments)              & 294 (3896) & 251 (2012) \\
            Trapped deaths, length 10+                     & 122        & 68         \\
            \bottomrule
      \end{tabular}
\end{center}
Free moves almost double the food. They also make our workers far longer at the unit cap: 40--140 where the old build's are
under 50. Long workers are juicy targets and die whole when trapped. Over full games, self-deaths of dragons of 16+ ran about
2.5$\times$ the old build's (137--169 vs 62--67 per game on big maps).

\subsection{Other checks}
\begin{itemize}
      \item The original scenario tests (\code{scen.py}) all give the same outcomes as before.
      \item Coil stress test (\code{stress_box.py}, 30 trials): no king self-deaths, for either build.
      \item Turn time (seed 20, 250 rounds, debug build): at most 6.1\,ms, 1.4\,ms on average (the Part I build: 2.7\,ms, 0.9\,ms).
      \item No crashes, timeouts or pay deaths in any game.
\end{itemize}

\subsection{How Part II got here}\label{sec:iter2}
\begin{enumerate}
      \item The first version let the king make long free moves to the window's edge. The two king deaths above led to
            \code{unseen_risk} and the estimate-based soft reach.
      \item \code{std::unordered_map<int, std::pair<int, Direction>>::operator[]} does not compile (\code{Direction} has no
            default constructor), so \code{insert_or_assign} and \code{at} are used.
      \item The route-end search first counted the other dragons' moves one turn ahead (offset $\code{turns_for}(t)$ instead of
            $\code{turns_for}(t) - 1$). That was slightly optimistic and is now fixed.
\end{enumerate}

%=====================================================================
\part{Landing safety and teammate memory}
%=====================================================================
\section{What was actually happening}
The concern was that long snakes, able to take 20 free steps, move far past their window and so move blindly into
teammates. They do not.
\begin{itemize}
      \item Free moves are capped at \code{FREE_MOVE_MAX_STEPS} $=6$ steps.
      \item Every step must be on a tile visible this turn, and \code{simulate_path} checks it against the board as it is.
      \item Dragons act one at a time, so nothing moves while we move.
\end{itemize}
What \emph{is} blind is the landing. A multi-step move can end at the window's edge. Before our next turn, whatever is
just outside the window (heads we cannot see, bodies we have forgotten) can close in, and the survival search cannot
see it either. Bodies were also forgotten after 2 rounds, however long they were.

\textbf{Measured} (\code{probe_freemove.py}, seed 20, 200 rounds): how often a dragon was boxed in (no legal move) on the
turn right after a move the survival search had called safe.
\begin{center}
      \begin{tabular}{@{}lrrr@{}}
            \toprule
            Length   & After a multi-step move & After a single step & Multi-step, after Part III \\
            \midrule
            10--29   & 28 / 2379 = 1.18\%      & 1 / 959 = 0.10\%    & 14 / 2536 = 0.55\%         \\
            30+      & 4 / 803 = 0.50\%        & 1 / 807 = 0.12\%    & 4 / 897 = 0.45\%           \\
            under 10 & 6 / 1673 = 0.36\%       & 18 / 4561 = 0.39\%  & 10 / 1592 = 0.63\%         \\
            \bottomrule
      \end{tabular}
\end{center}
The walls around those heads were allies' bodies on 569 sides, our own body on 451, kelp on 123, ally heads on 81 and
enemies on 44. Direct collisions with allies were rare (about 5 per 200 rounds).

\section{Body memory}
\begin{itemize}
      \item \code{CellMemory::busy_until}: the round from which a body seen on the tile should be gone. It is set to
            $\mathit{now} + \code{OCCUPIED_MEMORY} + 1$ when the body is seen, which is the old 2-round memory.
      \item \code{remember_bodies()}, after \code{observe_heads}, refines it. Segment $i$ of a visible head's body lasts
            $\lceil(\mathit{len} - i)/\mathit{pace}\rceil + 1$ rounds, at most \code{BODY_MEMORY_MAX} $=30$, where:
            \begin{itemize}
                  \item $\mathit{len}$ is \code{known_length} (for an ally, its last status report within 10 rounds, if longer);
                  \item $\mathit{pace}$ is half the dragon's free steps for an ally (ours take several steps on most turns),
                        and 1 for anyone else.
            \end{itemize}
      \item A visible segment of an ally whose head we cannot see, but whose length we know, lasts half that length at
            its pace, clamped to 3--\code{UNPLACED_BODY_MEMORY_MAX} $=8$ rounds.
      \item \code{body_remembered(m)} and \code{body_rounds_left(c)} replace the old fixed 2-round test. They are used by
            \code{remembered_occupied}, \code{blocks_pearl_path}, the time-aware space search and blind portal exits.
      \item The survival search frees a remembered tile after its rounds left (was a fixed 3); routing's \code{body_free_at_} does the same.
      \item The crowd count around route pearls still uses recent sightings only.
\end{itemize}
A first version assumed every body moves one step a turn. Measured over three maps, that was \emph{worse} than having no
long memory: our allies move faster, so it kept tiles closed too long and made the survival search pessimistic.
The paced version measured slightly better than no memory.

\section{Where allies can be before our next move}
In \code{prepare_survival}, the step-1 head reach now uses what we know:
\begin{itemize}
      \item \textbf{Allies}: $\min(\code{free_moves}(\code{known_length}), \code{FREE_MOVE_MAX_STEPS})$, since our allies never
            move further. \textbf{Enemies}: \code{free_moves(est)}.
      \item \textbf{Traced bodies} of visible heads use \code{known_length}.
      \item \textbf{Allies out of sight} that sent a status packet this round or the last: every tile within their reach of the
            reported head is blocked for our first step. The reach is doubled when the packet is from last round and the ally
            moves before us, since it has moved once since. The reported tile itself holds its body for
            $\min(30, \mathit{len} - 1)$ moves.
      \item \code{tiles_within} gained \code{any_first} (direction unknown) and now uses a set.
\end{itemize}

\section{Landing next to tiles we cannot see}
\code{unseen_risk} is now charged for every dragon from \code{position_value}, in two parts. With $n$ the tiles within
\code{UNSEEN_RADIUS} $=3$ steps of the end that we cannot see now:
\begin{itemize}
      \item \textbf{Ram} (valuable dragons, as in Part II): $\min(0.3,\ n \cdot \code{P_UNSEEN_ATTACK} \cdot \mathit{enemy\ density}) \times \code{death_cost()}$.
      \item \textbf{Box-in} (any dragon of \code{UNSEEN_BOX_MIN_LENGTH} $=10$+):
            $\min(0.3,\ n \cdot \code{P_UNSEEN_BOX} \cdot (\mathit{ally} + \mathit{enemy\ density})) \times \code{trap_cost()}$,
            with \code{P_UNSEEN_BOX} $=0.05$. \code{ally_density_x1e4_} is a running average like the enemy one.
\end{itemize}
Every single step has $n = 1$, so single steps are not affected relative to each other.

\textbf{Tried and dropped}: a hard rule that long dragons end a free move at most 2 tiles from where they started.
It measured no better than the price above.

\section{The late-game king keeps to its feeders' pace}
From \code{GATHER_ROUND} (380), the king and growers take no free moves while their goal is EXPLORE (outside a fight).
They still take them to eat along a route, to fight, and to escape. Without this, the \code{feeder_near} scenario failed:
a king of length 15 roamed 3--4 steps a turn and its length-4 feeder, at 1 step a turn, never reached it.

\section{Part III results}
\textbf{Six large maps}, 200 rounds each, both builds against the Part I build (\code{ab_trap.py}, seeds 20, 30, 37, 23, 27, 50):
\begin{center}
      \begin{tabular}{@{}lrrr@{}}
            \toprule
                                           & Part II build & Part III build & Change  \\
            \midrule
            Trapped deaths, length 10+     & 581           & 497            & $-14\%$ \\
            Accidental self-death segments & 18054         & 16669          & $-8\%$  \\
            Total length at round 200      & 5237          & 5334           & $+2\%$  \\
            Segments lost to enemy rams    & 14320         & 15664          & $+9\%$  \\
            \bottomrule
      \end{tabular}
\end{center}
More long snakes stay alive, and some of them are then rammed instead.

\textbf{Head-to-head} against the Part II build, full games on large maps:
\begin{itemize}
      \item the first version of these changes: 2--2;
      \item the paced memory: 1--3, every game by a small margin;
      \item the final build, seed 20 both sides: 1--1.
\end{itemize}
This is about even: the gain shows up in long-snake survival, not yet in wins.

\textbf{Also found}: every ``ran into another body'' death of ours is a single step through a portal whose exit we
cannot see, onto an ally's body (\code{probe_hitbody.py}). About 40\% of blind crossings die, the same before and after
this round. This risk predates v21 and is priced at \code{P_BLIND_PORTAL} $=0.3$.

\textbf{Checks}: the scenario tests give the same outcomes as before, including \code{feeder_near}, where the king is
fed again. The cleaned-up build was confirmed move-for-move identical to the tested one before the late-game change.

%=====================================================================
\part{Going back for remembered pearls}
%=====================================================================
\section{What was wrong}
In a game, a long king knew of pearls behind it but never went back for them.
\begin{itemize}
      \item \textbf{Range in steps.} Routes considered only pearls within \code{KEEPER_ROUTE_RADIUS} $=20$ steps (whole route:
            28). A king of length 80 takes 20 free steps a turn, so a remembered pearl 30--40 steps away around its own
            body was never a route. The king fell back to exploring, and nearby unexplored ground (cost 12) always looked
            closer than a pearl 30+ steps away.
      \item \textbf{Speed capped by the window.} Even with a route, a free move only crosses tiles visible this turn.
            Towards a target outside the window that is 3 tiles of progress a turn, whatever the free steps.
\end{itemize}
\textbf{Test} (\code{scen_route.py loopback}): a king of length 40 is made to walk 60 steps west past three pearls,
leaving them 22--25 steps behind with its body between, and then plays on its own.
\begin{center}
      \begin{tabular}{@{}lp{9.6cm}@{}}
            \toprule
            Build            & What happens                                                                                  \\
            \midrule
            Before           & No route: it explores away south and never eats them (20 turns watched).                      \\
            Route range only & Goes back at 3 tiles a turn and eats all three 8 turns after taking over.                     \\
            Final            & Goes back and rushes 8 steps at once when the ground is quiet: all three eaten in 6--7 turns. \\
            \bottomrule
      \end{tabular}
\end{center}
In this test the rush waits a few turns because the measured averages start from crowded-game values (see below).

\section{Route range in turns}
\code{route_radius()} and \code{route_horizon()} replace the fixed values: the old ones, or this many turns of our free
steps, whichever is further:
\begin{itemize}
      \item keepers: 4 turns (radius) and 6 turns (horizon);
      \item workers: 3 and 4 turns;
      \item capped at \code{ROUTE_MAX_RADIUS} $=100$ and \code{ROUTE_MAX_HORIZON} $=140$ steps.
\end{itemize}
A king of 40 now plans 40 steps ahead (60 for the whole route); one of 80, 80 and 120. Our time-aware map
(\code{route_head_dist_}) is searched out to the same radius. Dragons with up to 5 free steps are unchanged.

\section{Rush}
\code{add_route_rush} runs when the goal is EAT. It follows the checked route path (\code{route_path_}) for all our free
steps, never through a portal.
\begin{itemize}
      \item \textbf{Moves weighed}: all of it, stopping on each route pearl, and the part we can see. A fully visible move of
            up to 6 steps is skipped, since it is already among the free moves.
      \item \textbf{Remembered tiles}: with \code{remembered_ok}, \code{simulate_path} also accepts a step onto a tile we cannot see
            now, if we have seen it, know the edge into it, and remember no body on it. Those tiles are returned as \emph{blind},
            and a pearl remembered on one counts as eaten (its tail staying put is the cautious case for self-collision).
      \item \textbf{Scoring}: \code{evaluate_path}, plus remembered pearls on blind tiles at their freshness, minus the chance of
            hitting something times \code{death_cost()}.
      \item \textbf{Safety}: a rush with more than \code{RUSH_MAX_RISK} $=8\%$ (valuable dragons: $2\%$) chance of a hit is not safe.
            Rushes are exempt from the known-space penalty, since they price unseen tiles themselves.
      \item The chosen move logs as \code{RUSH}.
\end{itemize}
\textbf{Chance a blind tile is taken}: a tile last seen $k$ rounds ago risks
\[ \min(k \cdot \mathit{fill\ rate},\ \mathit{occupancy},\ \code{P_BLIND_TILE_CAP} = 0.5), \]
and a tile our own body is on now risks 0. Nobody else can be there, and nothing else moves while we do, so following our
own body round is safe however far it goes. Both rates are measured in \code{remember_vision} as running averages:
\begin{itemize}
      \item the fill rate: of the tiles we saw free last round and see again, the share now taken by another dragon;
      \item occupancy: the share of tiles in view holding another dragon.
\end{itemize}
They start at 3\% and 20\%. In a real game around the king they ran 3\%/round and 23\% in the midgame, and 1.5\% and 15\%
late. So in crowds a tile seen 7+ rounds ago is a 20\%+ risk and rushes are rare; in quiet ground or along our own body
they go.

\section{Bodies running out of sight}
\code{predict_hidden_segments()} runs after \code{remember_bodies}.
\begin{itemize}
      \item A visible segment's facing points to the next segment towards its head. If that tile is out of sight, it holds the
            same dragon.
      \item If no visible segment follows a segment towards its tail, and only one tile out of sight could, that tile almost
            surely does.
      \item Both tiles are remembered as taken for as long as the segment next to them, so rushes, the survival search and
            routing treat them as bodies. A remembered body no longer needs the tile to have been seen.
\end{itemize}
A rush in testing hit an ally's segment 3 just outside the window, where its body ran out of sight. Rush collisions went
from 5 to 2 in 300 rounds with this.

\section{The 12-bit id bug (predates v21)}
Status packets, beacons and king records carry dragon ids in 12 bits. Both teams share one id counter, and with free
moves (more eating, more splitting) it passed 4800 in a game. The king, id 4821, was reported as 725. Feeders looked
for a head with id 725, never found it, never fed it, and wandered at lengths 150--215 until they died in their own coils.
The team was wiped out in 2 of 6 games.

The fix: \code{pid(id)} $=$ \code{id & 4095} and \code{my_id()}. Every id we store (\code{SeenHead}, signals, the king
record, the ally length cache) and every comparison between a stored id and a real one uses that form; comparisons
between two real ids (a visible part against our own id) are unchanged. Two live dragons whose ids differ by exactly
4096 would be confused, which is rare.

\section{Part IV results}
Full games, large maps, both sides:
\begin{center}
      \begin{tabular}{@{}>{\raggedright}p{4.2cm}>{\raggedright}p{2.6cm}>{\raggedright}p{4.6cm}p{3cm}@{}}
            \toprule
            New build                              & Against        & Result                                   & Final max length (avg) \\
            \midrule
            Range, rush, hidden bodies (no id fix) & Part III build & 4--2, both losses team wipeouts (id bug) & --                     \\
            Final (with the id fix)                & Part III build & \textbf{4--0} (seeds 30, 37)             & 413 vs 245             \\
            Final                                  & Id fix only    & \textbf{5--1} (seeds 20, 30, 37)         & 348 vs 278             \\
            \bottomrule
      \end{tabular}
\end{center}
So both parts pay: the id fix keeps the king fed, and the longer routes and rushes add about a quarter to its final
length on top.

Over six maps of 200 rounds (before kings matter) the variants differ by about $\pm 3\%$ in total length, within
run-to-run noise. Rushes eat about 3\% more pearls there.

\textbf{Checks}: the scenario tests are unchanged (\code{room}, \code{contest}, the feeder tests) and \code{loopback} passes.
The cleaned-up build plays move-for-move like the tested one.

%=====================================================================
\part{The king's own moves}
%=====================================================================
\section{Deep self-safety}
Until now a move was ``proven safe'' when the survival search found $\min(L+2, 20)$ moves. For a king of 200 that proves
almost nothing: its newest segments stay put for up to 200 moves, so it could walk into a ring of its own recent body
holding 25 tiles and be called safe.
\begin{itemize}
      \item \code{deep_dragon()}: the king (also before growers start at round 300) and growers.
      \item For them, \code{survival_from(..., deep)} searches to $\min(L + 2, \code{KING_DFS_MAX_DEPTH} = 300)$. The budgets are
            \code{KING_DFS_NODE_BUDGET} $=20000$ per search and \code{KING_DFS_TURN_BUDGET} $=250000$ per turn
            (\code{KING_ROUTE_DFS_TURN_BUDGET} $=80000$ for route ends, which are deep for them too).
      \item In deep searches, another dragon's body that we cannot place (head out of sight) is assumed gone after
            \code{DEEP_UNPLACED_BODY} $=40$ moves rather than never, or a long search would fail on it.
      \item \code{survival_from} now takes a budget pool (0 moves, 1 route ends, 2 plans) and a per-search node cap.
\end{itemize}
\textbf{Test} (\code{scen_route.py ownpocket}): a king of 60 walks (scripted, so it knows its body) east and then round a full
ring; at hand-over its newest 23 segments enclose 25 tiles holding pearls, and its head is at the ring's one gap.
\begin{center}
      \begin{tabular}{@{}lp{10cm}@{}}
            \toprule
            Build         & Stepping into the ring                                              \\
            \midrule
            Part IV build & \emph{safe} (20 of 20 moves proven); kept out only by its score.    \\
            Final         & \emph{unsafe}: only 23 of the 62 moves needed exist; heavy penalty. \\
            \bottomrule
      \end{tabular}
\end{center}
Both builds wait outside and eat the pearls once the ring has moved on, with 3 pearls inside or 7.

\section{A body we do not know}
A dragon born long from a split (an heir) starts as a new process: no map memory, and of its body only what it can
see. On average heirs knew 79\% of their body (\code{probe_bodyk.py}). The rest is attached where the known body ends
(\code{own_trail_.back()}). In deep searches, while some of it is still there, a tile we have never seen within
$\min(\mathit{unknown}, \code{UNKNOWN_BODY_RADIUS} = 10)$ steps of that end counts as possibly our own body. (Multi-step
moves never caused a body-model reset; \code{probe_bodyk.py} counted none.)

\section{The king's plan}
\code{plan_king_moves} runs after all move candidates exist, for deep dragons with 2+ free steps.
\begin{itemize}
      \item For each of the \code{PLAN_CANDIDATES} $=8$ best moves (safe ones first), the turns after it are searched.
            For turns 2..\code{PLAN_TURNS} $=3$, a state's options are a shortest path to every tile within our free steps
            (at most \code{PLAN_MAX_STEPS} $=24$, no portals). The search is time-aware against our own moving body. The
            \code{PLAN_BRANCH} $=4$ best are kept, by food on the way and progress along the route (\code{plan_options}).
      \item Each option is played out exactly by \code{plan_extend}. Every step is checked against our body as the step before
            left it, tail included. A pearl, seen or remembered, eaten keeps the tail.
      \item Other bodies are timed cautiously (\code{plan_free_at_}): a tile another head can reach may hold its body for its
            length, not just for our next move. A plan should not count on gaps another dragon can fill.
      \item \textbf{A plan counts only if it is proven safe}: from where it ends, the deep survival search (budget pool 2,
            \code{PLAN_DFS_TURN_BUDGET} $=150000$) must find a way on for our whole length.
      \item The best safe plan's \emph{food} in its later turns, discounted \code{PLAN_DISCOUNT} $=0.85$ a turn, is added to the
            move (\code{Candidate::plan_value}, $\geq 0$; $-1$ means none was safe). Progress is left to the move's own score.
\end{itemize}

\textbf{Two dead ends on the way}
\begin{itemize}
      \item \textbf{Counting progress in the plan made the king hold back.} Later turns ``made'' 10 steps of progress each through
            tiles the real move would price as risky, so a single step now left more for the plan to collect. In
            \code{loopback} the king crept one step a turn (pearls eaten at turn 78, not 66). Food only fixed it.
      \item \textbf{Cautious timing in every deep search} (not just in plans) made kings doubt every move in crowds. It lost
            3 of 3 games against the Part IV build and was dropped. Used only inside plans it went 4--2; without it, plans
            went 1--5.
\end{itemize}

\section{How kings actually get trapped}\label{sec:cascade}
\begin{itemize}
      \item \textbf{Self-deaths of long dragons by role} (\code{probe_selfdeath.py}): about 9 per game are kings or growers.
            Most are dragons born from a split one or two turns earlier, already boxed in.
      \item \textbf{Escape splits} (\code{probe_kingtrap.py}, \code{probe_ksplit.py}): the king escape-splits on about
            \textbf{1 in 10} of its turns (56--89 times a game, in both builds). Each split costs the head and 2 segments and
            restarts the king as a new process that knows little of its body.
      \item \textbf{Where they come from}: about two thirds are ``blind'' splits by young heirs that do not know their body.
            These are chains: the head alternates between the two ends of a body boxed in at both. The rest are kings that
            knew their body and were boxed in mostly by \emph{other} dragons, often right after a multi-step move
            (\code{probe_kroom.py}).
      \item The deep check, the unknown-body caution and the plan did not change that rate.
\end{itemize}

\section{Part V results}
\begin{itemize}
      \item \textbf{Final build vs the Part IV build}, six full games on large maps, both sides: \textbf{4--2}. Final king length
            averaged 346 vs 338, which is within the large game-to-game swings.
      \item Scenario tests unchanged; \code{loopback} still eats at turns 66--67; \code{ownpocket} as above.
      \item Coil stress test: no king self-deaths.
      \item Turn time: at most 36\,ms (the king's plans and deep searches), 1.7\,ms on average.
\end{itemize}

%=====================================================================
\section{Known limitations and not yet done}
\begin{itemize}
      \item \textbf{Free-steps assumption}: per turn, from the length at the start of the turn. If the real engine counts differently,
            \code{free_moves} and \code{engine.py} are the two places to change.
      \item \textbf{Queens / win condition}: not implemented.
      \item \textbf{Teammate memory is partial}: we remember what we have seen and what status packets tell us (head position and
            length only). Where an unseen ally's body lies is guessed.
      \item \textbf{Blind portal crossings} die about 40\% of the time in crowded games.
      \item \textbf{Rushes} use only free steps. Paying segments to rush to a remembered pearl is not considered: the pearl may be gone.
      \item The fill rate and occupancy start at crowded-game values, so a newborn long dragon does not rush for its first ~30 turns.
      \item \textbf{The king's escape-split chains} (\S\ref{sec:cascade}): about 1 split in 10 king turns, mostly by heirs
            that do not know their body. The most valuable next step for the king.
      \item \textbf{Crowding at the unit cap} is now the biggest leak. Workers can no longer split there and grow very long with free moves;
            they then die whole when trapped or rammed. This is the most valuable next step (for example, feeding long workers to the king
            or capping worker length).
      \item The fight search still gives the enemy single-step replies plus a sprint into our head; free repositioning of 2+ steps is not played out.
      \item Free moves avoid portals and stay inside our window.
      \item \code{unseen_risk} counts unseen tiles; it does not use what we saw there last turn.
      \item Route ends: the other dragons' future moves are unknown. A body whose head we cannot see stays a wall forever.
      \item Edges never seen are assumed open (as before), so a route into unseen ground can look better than it is until we see it.
      \item The survival search covers our own single steps only; other dragons are bodies that drain away plus a head reach.
      \item Pearls in the survival search count only if seen this round.
      \item Planned but not started: panic fallbacks.
\end{itemize}


%=====================================================================
\part{v21 $\rightarrow$ v22: the judge's points budget}
%=====================================================================
\section{What the judge measures}
From the Timeouts page and the organisers' public repository (\code{unswcpmsoc/unswbc}, \code{pythoncli/unswbc/metering.py}
and \code{sandbox.py}, which reproduce the judge):
\begin{itemize}
      \item Bots run as WebAssembly (C++: \code{-O2 -msimd128}, libc++). Each instruction has a fixed price in \emph{points}: most cost 1,
            loads and stores 2, division 3, calls 4--6, memory growth 50, bulk copies 10 plus 1 per 8 bytes.
      \item \textbf{100 million points per dragon per turn}, no carry-over. The turn runs from reading the round to the flush in
            \code{end_turn()}; work after that counts for the next turn. A process's first turn includes its own setup.
      \item A write to stdout costs 2.5M plus 4,000 a byte; a read 6 a byte. The helper buffers the whole turn into one write (one
            2.5M charge a turn); \code{output_log} lines cost only their bytes.
      \item \textbf{The clock counts points}: \code{clock_time_get} returns points spent (1\,ns = 1 point). So
            \code{std::chrono::steady_clock} inside the bot is an exact meter of its own spending.
      \item A turn over budget is stopped and delivers nothing. The action stays at its default, suicide, and the engine records
            \code{DragonDeathReason::NoValidAction}: the ``no valid action'' deaths.
\end{itemize}
The engine (\code{engine/src/actions.cc}) also uses \code{NoValidAction} for a paid step at length 2, a split that would leave
fewer than 2 segments on either side, and a split at the unit cap. v21 already avoids all three: paid steps are capped at $L-2$
(\code{max_turn_steps}), and every split goes through \code{can_split}, whose unit count is current at our turn. So the deaths
were the points budget.

\section{Measuring it}
\texttt{unswbc run -{}-sandbox} measures exactly, but the toolkit ships through PyPI, which this workspace cannot reach. So the
judge's build was reproduced:
\begin{itemize}
      \item wasi-libc, libc++ and compiler-rt built from source for \code{wasm32-wasi} with clang 18 (\code{wasm/setup_toolchain.sh});
            the bot built with \code{-O2 -msimd128 -mbulk-memory} (\code{wasm/build_bot.sh}). Exceptions are off; the judge's
            clang 20 with wasm exceptions may cost a little more.
      \item The organisers' own \code{metering.py} adds the point counter (\code{wasm/meter/instrument.py}).
      \item \code{wasm/meter/runner.mjs} runs it under Node's WASI the judge's way: the clock returns points, writes and reads are
            charged, and each turn's points are logged.
      \item The bot is deterministic, so a game is recorded natively (each dragon's input, \code{rec_bot.sh}) and every dragon is
            replayed through the meter (\code{replay.py}, \code{analyze.py}). A \code{BOT_PROFILE} build adds the points of
            each phase (\code{prof_an.py}).
\end{itemize}

\section{v21's turns}
One recorded 500-round game (seed 20, 56$\times$53): our dragons of length 30+ and a sample of shorter ones, 7,410 turns,
in millions of points.
\begin{center}
      \begin{tabular}{@{}lrrrrrr@{}}
            \toprule
            Length   & Turns & Median & p90   & p99    & Max    & Over 100M \\
            \midrule
            2--9     & 1715  & 20.3   & 33.3  & 81.9   & 150.6  & 13        \\
            10--29   & 1700  & 28.6   & 53.0  & 95.5   & 116.3  & 8         \\
            30--59   & 2702  & 40.8   & 102.7 & 173.3  & 299.2  & 292       \\
            60--99   & 933   & 39.6   & 200.4 & 351.5  & 934.7  & 252       \\
            100--149 & 210   & 27.7   & 312.0 & 536.4  & 840.8  & 75        \\
            150--199 & 77    & 30.4   & 575.4 & 1217.5 & 1217.5 & 31        \\
            200+     & 73    & 285.6  & 578.8 & 850.8  & 850.8  & 61        \\
            \bottomrule
      \end{tabular}
\end{center}
Kings averaged \textbf{164M a turn} (max 1.2B). Of that:
\begin{itemize}
      \item 84M (51\%) was the king's plan, almost all in \code{plan_options}: for every tile within 24 steps it copied the path
            and looked for the tile along the whole body (up to 300 segments) and the route.
      \item 50M (31\%) was routing: 14 whole-map searches for the pearls' fields and the route replays.
      \item A survival-search position cost about 2,800 points, mostly \code{step_from} (edge copies, portal lookups in a
            \code{std::map}).
\end{itemize}
Short dragons sat at 15--20M a turn, most of it map-wide setup.

\section{The turn's clock and deadlines}
\begin{itemize}
      \item \code{clock_now()} is \code{steady_clock} in nanoseconds: points in the judge. Native builds (local tests) multiply by
            \code{NATIVE_POINTS_PER_NS} $=24$, measured against the metered build (p10--p90: 17--32).
      \item \code{main} starts the turn's clock before reading the round (the first turn: at program start).
      \item Deadlines, counted from the turn's start:
            \begin{itemize}
                  \item the route's checks stop at \code{ROUTE_STOP} $=40$M;
                  \item the moves' survival searches at \code{MOVES_STOP} $=66$M (first 64M);
                  \item the king's plans at \code{PLAN_STOP} $=91$M (first 82M);
                  \item anything else at \code{SEARCH_STOP} $=92$M (first 86M).
            \end{itemize}
            The results below were measured with the first values. \S\ref{sec:peak95} raised them to a 95M peak.
            The phases are absolute, so a turn whose setup is expensive just leaves less for the later phases.
      \item One survival search may spend at most a quarter of what its phase has left (at least \code{SEARCH_MIN_POINTS} $=2$M).
            A search deep in a tight spot can no longer use up the phase before the escape sprints are searched.
      \item Where the clock is read:
            \begin{itemize}
                  \item \code{survival_from}: before starting, and every 128 positions;
                  \item \code{plan_king_moves}: before each candidate and each state's options;
                  \item \code{plan_route}: after each pearl's field (fewer pearls are kept when time is short) and before each replay.
            \end{itemize}
      \item \textbf{A search cut by time is treated like one cut by its node budget}:
            \begin{itemize}
                  \item a candidate's survival falls back to the room estimates;
                  \item a route end falls back to the free-neighbour check;
                  \item a plan that cannot be proven adds nothing.
            \end{itemize}
            The node budgets stay; the deadlines only make sure the turn is never over.
\end{itemize}

\section{Cheaper searches}
\begin{itemize}
      \item \textbf{Neighbour table.} \code{step_from} for every tile and direction lives in \code{nbr_} (the tile stepped onto,
            \code{NBR_WALL}, or \code{NBR_BLIND} for a portal whose far end we do not know), plus \code{nbr_portal_}. It is kept
            from turn to turn and refreshed once a turn, after the sonar and the vision are learned.
            \begin{itemize}
                  \item Only tiles whose edges changed are rebuilt, with their four neighbours: a step reads the far side's edge when
                        its own is unknown.
                  \item A newly learned portal end rebuilds everything.
                  \item \code{step_from} itself reads the table, so every caller got faster. The survival search, the plan, the route
                        fields, \code{timed_bfs}, \code{head_eta_bfs} and \code{replay_leg} index it directly.
                  \item \code{NBR_CHECK} compares it against the uncached \code{step_from} after the refresh and again mid-turn: no
                        mismatch in two games. A deliberately broken refresh gives 133.
            \end{itemize}
      \item \textbf{\code{plan_options} in one pass.}
            \begin{itemize}
                  \item Our body's tiles go in a stamped table (O(1) per tile instead of a walk along the path).
                  \item Food and the route pearls eaten (a bitmask) are carried along the search tree.
                  \item The route's remaining length is cached per set of pearls eaten: only its first leg depends on the tile.
                  \item Only the best \code{PLAN_BRANCH} options are kept, without sorting all of them, and nothing is allocated per call.
                  \item One behaviour change: pearls on tiles the plan already ate no longer count as food again.
            \end{itemize}
      \item \textbf{Hash maps out of hot loops}: \code{replay_leg}'s parents and \code{head_eta_bfs}'s segment indices are
            stamped arrays, and \code{route_bfs} uses a vector queue.
      \item \textbf{The survival search's body is a ring buffer} instead of a \code{std::deque}. \code{SURV_CHECK} runs both on
            every search: no mismatch; a broken \code{back()} gives 22. It saved only about 4\% a position.
\end{itemize}
A survival-search position now costs about 1,100 points (was 2,800). For kings of length 60+, \code{plan_options} fell from
103M to 5--12M a turn and the route fields from 43M to 8--12M.

\section{Part VI results}
\textbf{Points} (the same game and dragons as above, release build):
\begin{center}
      \begin{tabular}{@{}lrrrrrr@{}}
            \toprule
            Length   & Turns & Median & p90  & p99  & Max  & Over 100M \\
            \midrule
            2--9     & 1715  & 14.4   & 18.5 & 31.4 & 47.5 & 0         \\
            10--29   & 1700  & 17.8   & 24.1 & 31.1 & 36.6 & 0         \\
            30--59   & 2702  & 21.5   & 33.4 & 43.0 & 83.5 & 0         \\
            60--99   & 933   & 22.0   & 44.9 & 64.7 & 85.4 & 0         \\
            100--149 & 210   & 21.3   & 60.6 & 85.5 & 85.8 & 0         \\
            150--199 & 77    & 23.1   & 85.3 & 86.0 & 86.0 & 0         \\
            200+     & 73    & 59.4   & 85.2 & 85.6 & 85.6 & 0         \\
            \bottomrule
      \end{tabular}
\end{center}
\begin{itemize}
      \item Kings now average 36.5M a turn: route 10.5M, plan 9.7M, free moves 4.5M, objective map 4.4M.
      \item Over four recorded games (15,225 turns of the dragons that reached length 40), the worst turn is \textbf{86.6M}.
            With the earlier 82M stop, every turn of every dragon in two v22 games (52,359 turns) stayed at or under 82.7M.
      \item Peak memory is 2.1\,MB of the 48\,MB allowed.
      \item The work after the last deadline (choosing, the output) is under 1M. The 13M margin is for the judge's build, which
            is not exactly ours.
\end{itemize}

\textbf{How much lookahead the budget costs.} On recorded inputs, v22 with deadlines off (\code{BOT_UNLIMITED}) chooses
the same move as v21 on 2,244 of 2,258 turns. The 14 differences are all in one 200+ king and come from the plan's new
tie-breaks. With deadlines on, v22's move differs from the unlimited one on:
\begin{itemize}
      \item 0 of 2,047 turns at length $<100$;
      \item 1 of 88 at 100--149;
      \item 7 of 83 at 150+.
\end{itemize}
At length 200+, the plan is cut on 8 of 73 turns and a single search's share on 14. The route and the moves are rarely cut
(3 and 1).

\textbf{Games.} \texttt{engine.py -{}-tle-ms} kills a dragon whose turn takes longer than that many wall-clock milliseconds, as the
judge would. 4.5\,ms is about 100M points at 24 points/ns. Each game ran on its own (\code{ab_tle.sh}).
\begin{center}
      \begin{tabular}{@{}>{\raggedright}p{4.6cm}>{\raggedright}p{1.3cm}>{\raggedright}p{3.6cm}p{5.4cm}@{}}
            \toprule
            Games                                                                      & Result        & Longest at the end (v22 / other) & Limit deaths per game \\
            \midrule
            v22 vs v21, limit on: seeds 20, 6, 27, 30, 20 (v22 second in the last two) & \textbf{5--0} & 81--139 / 0 (wiped out)          &
            v22: 1--2. v21: 105--278, of them kings 49--99                                                                                                        \\
            v22 vs v22, limit on: 2 games                                              & --            & --                               & 0                     \\
            v22 vs v21, \emph{no} limit (v21 unmetered): 3 games                       & 0--3          & 170--194 / 392--456              & --                    \\
            v22 with deadlines off vs v21, no limit: 2 games                           & 1--1          & 251, 354 / 271, 299              & --                    \\
            \bottomrule
      \end{tabular}
\end{center}
Two more limited games, recorded for the replays above (seeds 6 and 27), also went to v22 with v21 wiped out.
\textbf{Correction (see \S\ref{sec:nativeclock}):} native v22 in these simulator games ran without its searches, because
of a bug in its native turn clock. So the ``no limit'' rows compare v21 with a v22 that had no lookahead, and do not measure
what the budget costs. The metered replays, the move comparisons (run from recorded files) and your sandbox runs were not
affected.

Under the limit, v21 is wiped out in every game. Its long dragons and kings die to the meter, which is what was seen in
real matches. Without a limit, v21's unmetered lookahead (up to 1.2B points on a long king's turn) does buy strength. v22
gets as much of it as 100M allows.

v22's own limit deaths (7 in the 5 games against v21, none in self-play) took 4.9--5.8\,ms, twice about 30\,ms, and one was not
timed. They are wall-clock noise of the local machine: native v22 stops its searches at about 3.6\,ms, and the metered
replays show no turn over 86.6M.

\section{Checked in the sandbox, and a 95M peak}\label{sec:peak95}
\textbf{The official sandbox} (\code{unswbc run --sandbox}, v21 as team A, v22 as team B):
\begin{center}
      \begin{tabular}{@{}lrrrr@{}}
            \toprule
            Team & p50   & p99   & Mean  & Max   \\
            \midrule
            v21  & 26.3M & 52.1M & 27.0M & 99.7M \\
            v22  & 17.1M & 25.8M & 17.1M & 85.3M \\
            \bottomrule
      \end{tabular}
\end{center}
\begin{itemize}
      \item v22 won easily: its longest dragon reached 72, v21's 34.
      \item Seven v21 dragons died to the limit, all of them kings. A turn over budget shows as just under 100M because the
            sandbox stops it there.
      \item v22's 85.3M peak agrees with the metered replays here (86.6M), so the rebuilt toolchain prices the bot like the judge.
\end{itemize}

\textbf{The peak is the plan's deadline plus a fixed tail.} The metered replays were re-paired line by line: on replays from a
file, the runner's round labels can run ahead, because the bot's input buffer reads ahead.
\begin{itemize}
      \item The plan overshoots its stop by at most 0.8M (it ended at 82.81M at most against an 82M stop).
      \item Everything after it costs at most 4.1M on the 73 turns whose plan ended past 75M. That is the
            choosing, \code{send_status}, and the 2.5M write plus 4,000 a byte. The profiling build's own line is included;
            the release build is about 0.7M less.
\end{itemize}
So \code{PLAN_STOP} $=91$M with \code{SEARCH_STOP} $=92$M. \code{SEARCH_STOP} is kept just above the plan's stop, because a
search after the plan (a split escape, say) could otherwise run to it and then pay the tail.

\textbf{Measured} over the same four recorded games (15,225 turns of the dragons that reached length 40):
\begin{itemize}
      \item p50 18.2M, p99 66.8M, max \textbf{95.05M};
      \item one turn over 95M, none over 96M.
\end{itemize}
This leaves 5M for anything the judge prices differently from the rebuilt toolchain. Its own run measured
85.3M where the replays gave 86.6M.

\section{Tools added}
\begin{itemize}
      \item \code{wasm/setup_toolchain.sh}: builds the wasm32 sysroot (wasi-libc, libc++, compiler-rt builtins) with the system clang.
      \item \code{wasm/build_bot.sh SRC OUT.wasm [flags]}: the judge-like build.
      \item \code{wasm/meter/instrument.py} (with the organisers' \code{metering.py}): adds the point counter.
      \item \code{wasm/meter/runner.mjs}: the metered runner (\code{METER_BUDGET}, \code{METER_LOG}).
      \item \code{wasm/rec_bot.sh}, \code{wasm/replay.py}, \code{wasm/analyze.py}, \code{wasm/prof_an.py}: record, replay and summarise.
      \item \texttt{sim/engine.py -{}-tle-ms}; \code{sim/ab_tle.sh}.
      \item Build flags:
            \begin{itemize}
                  \item \code{BOT_PROFILE}: a \code{PROF} line of points per phase, sub-timers, nodes and the deadlines hit;
                  \item \code{NBR_CHECK} and \code{SURV_CHECK}: the equivalence checks above;
                  \item \code{BOT_UNLIMITED}: no deadlines, for tests only.
            \end{itemize}
\end{itemize}
With the official toolkit installed, \texttt{unswbc run \textless{}map\textgreater{} a b -{}-sandbox} gives the judge's own numbers. It is worth running once
on a long game to confirm the margin.

\section{v22: known limitations and next steps}
\begin{itemize}
      \item \textbf{Setup costs every dragon 12--15M a turn}: objective map, survival preparation, ETAs and parsing. It is
            unguarded but bounded by the map size; trimming it would give every phase more.
      \item Long kings (200+) still have their plan cut on about one turn in nine (8 of 73). The next savings are free moves (about 15M on
            those turns), the route fields (12M) and the plan's option search.
      \item The deadlines are fixed shares of one turn. Search depth could be traded between phases by role. The budget
            resets every turn; nothing can be saved for a hard turn.
      \item The native factor (24 points/ns) is a median. Native test games only approximate the judge; the metered replay is
            the reference.
\end{itemize}


%=====================================================================
\part{v22 (continued): the queen}
%=====================================================================
\section{The rule}
From the engine (\code{engine/src/scoring.cc}), at the last round the ranking is:
\begin{enumerate}
      \item the queen's length (dragon 0 or 1, the first dragon each team's map lists; a dead queen counts 0);
      \item the longest dragon;
      \item total length.
\end{enumerate}
A team that is wiped out loses at once. A split keeps the parent's id with the head part, so a queen that never splits stays the queen.
The official maps give each team 1 to 9 starting dragons (\code{arena} and \code{Colosseum}: 1; \code{default}: 4; \code{australia}: 9).

\section{What the queen does (\texttt{queen.hpp})}
\begin{itemize}
      \item \textbf{Who}: \code{is_queen()} is \code{controller_.get_id() < 2}, the raw id. Our packets' 12-bit ids wrap past 4096, so
            heads we see carry their own \code{SeenHead::queen} flag, and beacons carry a queen bit (bit 18).
      \item \textbf{No splits}: no worker split (\code{should_split}), escape split (\code{play_split_escape}) or pocket raid
            (\code{pocket_plan_allowed}). Trap pricing assumes she cannot split out (\code{trap_cost}).
      \item \textbf{Growing the swarm first} (revised; see \S\ref{sec:seed75}): until round \code{QUEEN_SEED_END_ROUND} $=75$ she plays
            as a founder and makes the usual small worker splits. The head part, and her id, stays with her. (First version: only
            when she started alone, until we were 4 or round 60.)
      \item \textbf{The cycle, from round 75 until \code{QUEEN_CYCLE_END_ROUND} $=400$}: a Hamiltonian cycle of a $w\times h$ block around her head
            (row 0, a serpentine over the other columns, back up column 0; $h$ even).
            \begin{itemize}
                  \item The block is the largest size from $16\times16$ down, at most $\max(144,\ 2(L + \code{QUEEN_CYCLE_SLACK}))$
                        tiles. Every step of its cycle must be an ordinary step (no kelp, no portal), by the neighbour table.
                  \item Of the placements, the one with the fewest other bodies on it and the fewest tiles never seen wins.
                  \item She walks it with as many of her free steps as are \emph{queen-safe}: the usual safety, plus no enemy, ally,
                        reach or soft threat at the end, and her whole length proven to have a way on (or two exits if that search
                        was cut). Paid steps are never used.
                  \item The route search and the king's plan are skipped while she cycles.
            \end{itemize}
      \item \textbf{Her body on the cycle}: walking it is safe only while every tile ahead is free of her own body by the time she
            reaches it (\code{queen_body_fits}). A segment $i$ places back is gone after $L - i$ of her steps, one more for each pearl she
            eats on the way (the known pearls ahead plus \code{QUEEN_GROWTH_MARGIN} $=2$); at least \code{QUEEN_CYCLE_MIN_FREE} $=4$
            spare tiles are required. She only follows the cycle while this holds. A new block (kelp found, or within 12 of her
            length) must pass the same test, and is never smaller than a cycle that still fits.
      \item \textbf{Off the cycle} (its next tiles are not queen-safe, or no cycle fits her): she takes the queen-safe move that ends
            nearest the cycle. If there is none, she takes the usual best move, as a deep-searching long dragon that never splits.
            She looks for a cycle that fits every \code{QUEEN_RETRY_ROUNDS} $=3$ rounds.
      \item \textbf{Crown, from \code{QUEEN_CROWN_ROUND} $=380$}: she stands for king, and she beats everyone in \code{offer_king} and
            \code{elected_king}, so beacons carry her and feeders gather on her. Before 380 she does not stand, and her length
            reports are ignored for the king/vice roles, so the team still has a roaming king. From 400 she retains as the king,
            everyone else feeds her, and her cycle ends. If she dies, the normal election takes over.
\end{itemize}

\section{A bug in the native turn clock}\label{sec:nativeclock}
The turn's clock started before reading the round. That is right in the judge: waiting costs no points there, and work after
\code{end_turn} is charged to the next turn.

Natively, though, the clock ran on while the bot waited for every other dragon to move. So in every native simulator game,
every turn began past all its deadlines and \textbf{v22 ran no searches}: no survival search, no plan, no route checks.
\begin{itemize}
      \item This was found when the queen's candidates all showed an unsearched survival result.
      \item Native builds now start the clock once the round is read (\code{main}). The WebAssembly build, the one the judge runs,
            is unchanged.
      \item Affected: the native games in Part VI (the 5--0 under the limit, the 0--3 without one, \code{loopback} being slower).
      \item Not affected: the metered replays (recorded files and the points clock), the move comparisons from files, and your
            sandbox runs.
\end{itemize}

\section{The simulator's pearls are far too rich}
On the official maps, pearl beds and respawn gaps are set tile by tile (\code{TILE x y minGap maxGap}). On \code{default.map}
most tiles respawn every 1--3,849 rounds (about 1,900 on average), and some every 1--255 or 1--385. The simulator's random maps
use 4--60 on almost every tile: 30--100$\times$ more pearls. That is why simulator dragons reach 300--450 while the official
sandbox game ended with a longest dragon of 72.
\begin{itemize}
      \item \code{engine.py --pearl-scale k} (and \code{probe_queen.py}'s 6th argument) stretch the simulator's gaps by $k$.
      \item The simulator now ranks by queen length, then longest, then total, and reports \code{A_queen}/\code{B_queen}.
      \item Everything tuned in the simulator so far was tuned for the rich economy. Loading the official \code{.map} files into the
            simulator is the proper fix.
\end{itemize}

\section{Queen results (simulator, so indicative only)}
\begin{itemize}
      \item With the simulator's rich pearls she grew to 35--77 by round 45--50, and every cycle filled up; she died off the cycle.
            This is what led to the body test above.
      \item With gaps stretched 30$\times$, she walks the cycle on most turns (48 of her first 50 in one game). Her deaths there were
            enemy rams at length 3--7 (rounds 51, 107, 189): an enemy queen's 2-step sprint, and workers' attack sprints, on turns
            when no move was queen-safe.
      \item Queen build against an identical build without the queen (\code{NO_QUEEN}), 4 founders, 30$\times$ gaps: 1--1. Side A was
            wiped out in both games whichever build it was (a quirk of the stretched simulator to look into). Its queen
            survived in the game it won.
      \item Metered (the judge's pricing), over 250 queen turns: 11.2M points on average, 37.1M at most.
\end{itemize}
These need checking on official maps in the sandbox, with one founder (\code{arena}) and several (\code{default}).


\section{The queen grows the swarm until round 75}\label{sec:seed75}
A non-splitting queen hurt our early game: on maps with few founders, one of them stopped growing the swarm. Now:
\begin{itemize}
      \item \textbf{Until round 75} she is an ordinary founder (\code{queen_seeding()} is just \code{is_queen()} and the round). She
            makes worker splits (she keeps her head, so she stays the queen) and joins the centre rush on small maps.
      \item She never makes pocket raids, and makes escape splits only when she is certainly dead (\S\ref{sec:doomed}).
      \item \textbf{From round 75}: her Hamiltonian cycle. From 380 feeders gather on her; from 400 she is the king.
      \item Gone: \code{queen_alone_} and \code{QUEEN_SEED_UNITS}.
\end{itemize}
\textbf{Simulator} (pearls 30$\times$ sparser; indicative only):
\begin{itemize}
      \item 32$\times$16, 3 founders, against the previous build: our swarm went 7 $\rightarrow$ 25 $\rightarrow$ 29 dragons (rounds 75,
            100, 125), while the other side, with a non-splitting queen, fell to 1. The game still went to the other side on queen length:
            ours died and theirs lived, which shows how much the first tiebreak weighs.
      \item 60$\times$46, 4 founders: she split until round 75, then walked her cycle on 86 of the next 100 turns, and was alive when the
            other team was wiped out.
\end{itemize}

\section{When the queen is certainly dead, she saves the rest of her length}\label{sec:doomed}
A split leaves the queen where she is, facing the same way: it cannot save her, only her length. So she splits only when
nothing she can do saves her (\code{play_split_escape}). Every candidate move must be one of:
\begin{itemize}
      \item a head trade;
      \item a move that leaves her head no free side (unless it eats a pearl, so the split can come next turn with one more segment);
      \item a move from which the survival search, run to completion, finds at most one more move.
\end{itemize}
Unlike other dragons, an enemy that could reach her does not count as certain death, since it may not come. When she is doomed:
\begin{itemize}
      \item the rear $L-2$ segments become a new dragon, if its end is not boxed in too;
      \item she, the 2 left behind, dies as she would have anyway;
      \item the ``cheap way out'' room estimates cannot veto it, because her test already searched every way on.
\end{itemize}
Pocket raids stay off for her, and trap pricing still treats a trap as her death.

\textbf{Test} (a test build that starts her cycle at round 1, in the simulator's rich economy, where she gets boxed in):
\begin{itemize}
      \item at length 102 she split off 100 (the child had 120 tiles of room);
      \item at length 54 she split off 52 (72 tiles);
      \item each time her head had no room left, and she died the next turn.
\end{itemize}

\textbf{Found while testing: seeding is risky for her.} As a founder she splits at length 4 and keeps 2, and she joins fights like
any worker. In three rich-economy games she died at length 2, boxed in among enemies, by rounds 26--39, in this build and the
previous one alike.

%=====================================================================
\part{v22 (continued): the centre rush on small maps}
%=====================================================================
\section{Which maps}
Top teams head for the centre on Autarky (54$\times$18), Devil, Portals, Prisoners Dilemma and Tower Defense (32$\times$16),
Stripes (24$\times$12) and Weakhold (40$\times$15), but not on the larger maps. Area does not separate them: Trophy (625) and
Queen of Spades (875) are smaller than Autarky (972) and are not on the list. \textbf{The short side does}:
\begin{itemize}
      \item the listed maps' short sides are 12--18;
      \item every other mid-size or large map's is 24 or more;
      \item the tiny maps (Arena 11$\times$11, Small 16$\times$8, Colosseum and Default Small 16$\times$16) fall under the same rule.
\end{itemize}
So the rush is on when $\min(W,H) \leq \code{CENTRE_MAP_MAX_SHORT_SIDE} = 20$. The bot knows $W$ and $H$ from its first turn.

\textbf{Added: compact maps.} The centre also pays on Trophy (25$\times$25, whose centre beds are 3.8$\times$ richer than its edges).
Its short side is 25, the same as Queen of Spades (25$\times$35), which is not on the list. Area cannot separate them either: Autarky
(972) is on the list and larger than Queen of Spades (875). So a second condition was added: the rush is also on when
$\max(W,H) \leq \code{CENTRE_MAP_MAX_LONG_SIDE} = 30$.
\begin{itemize}
      \item \textbf{On}: Arena, Small, Colosseum, Default Small, Stripes, Devil, Portals, Prisoners Dilemma, Tower Defense, Weakhold,
            Trophy, Autarky.
      \item \textbf{Off}: Queen of Spades (both), Default (32$\times$32), Maze, Stronghold, Trauma, Slithery Fight, Islands, Schooltime,
            and the 64$\times$64 maps.
\end{itemize}

It is not about pearls. On Prisoners Dilemma, Weakhold and Autarky the centre's pearl beds give 2--9\% of the outer rate;
Tower Defense's give 15$\times$ it. On every listed map the teams start at opposite ends of the long axis, so the centre is the
ground between them.

\section{What it does}
\begin{itemize}
      \item For the first \code{CENTRE_RUSH_END_ROUND} $=75$ rounds, every dragon gets \code{CENTRE_PULL} $=300$ (thousandths of a
            segment) per step a move brings it closer to the centre. 450 was tried; 300 was kept.
      \item Exploring is worth 180 a step plus up to 120 for dispersal, so the centre wins over exploring. A pearl next to the
            dragon still wins over the centre.
      \item ``Closer'' is by path: a search from the centre over the neighbour table each turn of the rush (\code{prepare_centre}),
            with kelp and portals as we know them and unseen edges open. A straight-line pull would press dragons against walls.
      \item The pull stops within \code{CENTRE_RADIUS} steps of the centre (short side / 4, at least 2), so dragons do not pile onto
            one tile. It applies to every goal but fights and feeding (\code{goal_progress} now adds it to
            \code{goal_progress_base}).
      \item After round 75 it is off, so nobody gets stuck in the middle.
      \item The queen walks her cycle and does not rush (her moves are her own; see the queen part).
\end{itemize}

\section{Simulator check (random maps, pearls stretched 30$\times$: indicative only)}
\code{probe_centre.py} gives the mean head distance to the centre (rush build vs current build):
\begin{itemize}
      \item 32$\times$16, 3 founders each: 4.0--5.7 vs 7.5--14.0 from round 25 to 75. Afterwards the rush build spreads out again
            (11.0 at round 100).
      \item 40$\times$15, 4 founders each: 4.6--7.2 vs 11.8--16.6 from round 10 to 75.
      \item Both games, and the second with the sides swapped, left the rush team with far fewer dragons by round 75 (4--5 against
            8--15). The likely reason is fewer pearls eaten on the way, so fewer splits. Whether this holds on the real maps, where top
            teams do rush, is for the sandbox to show; \code{CENTRE_PULL} is the knob to try lower (around 150, so pearls compete).
\end{itemize}


%=====================================================================
\part{v22 $\rightarrow$ v23: per-map settings and the queen's survival}
%=====================================================================
v23 starts here. Everything before this part is v22.

\section{A map table replaces the size rule}
Autarky (54$\times$18) turned out to be a near-certain loss with the centre rush, while Trophy (25$\times$25) gains from it. No size
rule fits that, and neither does a pearl one: Autarky's centre beds give 3\% of its outer rate, but Weakhold's (2\%) and Prisoners
Dilemma's (9\%) are just as poor, and the rush pays on both. So the choice is now made per map:
\begin{itemize}
      \item \code{Config::CENTRE_MAPS} lists the sizes with the rush: 11$\times$11 (Arena), 16$\times$8 (Small), 16$\times$16 (Colosseum,
            Default Small), 24$\times$12 (Stripes), 32$\times$16 (Devil, Portals, Prisoners Dilemma, Tower Defense), 40$\times$15
            (Weakhold), 25$\times$25 (Trophy). The bot knows the size from its first turn.
      \item Among the 25 official maps, maps of the same size all get the same answer, so the size identifies the map well enough.
            Should two of the same size ever need different settings, a founder can also tell them apart by its starting tile.
      \item \textbf{Any other size gets no rush}: Autarky, the larger maps, and any map we have never seen (the grand finals may use new ones).
      \item \code{CENTRE_MAP_TABLE_ENABLED} $=$ \code{false} turns the rush off on every map.
      \item Gone: \code{CENTRE_MAP_MAX_SHORT_SIDE} and \code{CENTRE_MAP_MAX_LONG_SIDE}.
\end{itemize}
Checked in the simulator:
\begin{itemize}
      \item 54$\times$18: the table build never came within 4 steps of the centre. The size-rule build reached a mean of 4.6 steps by
            round 50, with 4 dragons in the middle.
      \item 32$\times$16: both sides rushed as before.
\end{itemize}


\section{Walled in, or only contested?}
\textbf{What happened} (from two sandbox screenshots). A 2-long worker heading down had three moves:
\begin{itemize}
      \item \textbf{right}: into a tile with kelp above and to its right, whose only open side was the head of an ally worker heading left;
      \item \textbf{down}: into the tile that worker was about to enter;
      \item \textbf{left}: open water.
\end{itemize}
It went right. The ally moved on, so its neck sealed the last side, and our worker died on its next turn whatever happened.

\textbf{Why.} The survival search counts every tile another head could reach on its next move as taken for our next step. With
two heads nearby, that took every tile beyond left and beyond down. So all three moves came out ``searched completely, 0 of 4
moves''. Left and down were judged exactly as dead as the kelp corner, and a small difference in goal progress picked the corner.

\textbf{Change.} A move whose search says dead is searched again (\code{survival}) with only real walls: kelp, bodies as they drain
away, our own body, but not other heads' reach (\code{wall_free_at_}, built before those marks in \code{prepare_survival}).
\begin{itemize}
      \item \textbf{Still dead}: \code{SurvivalResult::walled}, copied to \code{Candidate::walled}. It is priced as certain death
            (\code{trap_cost()} in full, not 0.9$\times$ the shortfall), so it can never be preferred to a move that is only contested.
      \item \textbf{Through}: the move is \emph{contested}. It takes the second search's result (a way on), and the usual threat
            prices (adjacent heads, reach for valuable dragons, ally congestion) judge the contest.
      \item \textbf{Second search not run} (out of nodes or points): the first result stands, as before.
\end{itemize}
It runs only on moves the first search rejects, so it costs little. The queen's ``certainly dead'' test uses the same results, so a
merely contested queen no longer counts as doomed.

\textbf{Test} (\code{sim/scen_walled.py}: the screenshot position rebuilt tile for tile; our worker's candidates):
\begin{center}
      \begin{tabular}{@{}lll@{}}
            \toprule
            Move                & Before (search, score)    & After (search, score)   \\
            \midrule
            right (kelp corner) & 0 of 4, complete; $-5514$ & 0 of 4, walled; $-6014$ \\
            down                & 0 of 4, complete; $-5175$ & 4 of 4, safe; $-675$    \\
            left                & 0 of 4, complete; $-4995$ & 4 of 4, safe; $-495$    \\
            \bottomrule
      \end{tabular}
\end{center}
Before, the three were within about 500 points of each other. Here left happened to win; in the real game, right did. Now the
dead end trails by about 5,500. A full simulator game against the previous build: no crashes, no timeouts.


\section{Kill the enemy queen on sight}
Their queen's length is the first tiebreak. So, in \code{queen.hpp}, every dragon checks this before anything else it might do
(\code{play_queen_kill}, first in \code{play_turn}):
\begin{itemize}
      \item \textbf{Who}: any of us that sees the enemy queen (an enemy head with raw id 0 or 1) in its 7$\times$7 view, except our own
            queen, who would die with her.
      \item \textbf{How}: our head onto hers kills her (head on head: we die too). Dragons move one at a time, so where we see her head
            is where it is when we move. The shortest way there this turn goes through tiles we can see with no dragon on them, at most
            \code{max_turn_steps} (free steps, then paid ones, keeping 2 segments). It is replayed exactly (\code{simulate_path}).
            If one exists, that is the move, ignoring roles, trade values and rounds. That includes the old ``no aggression before
            round 50'' rule, and kings and growers.
      \item \textbf{Out of reach}: \code{QUEEN_HUNT_PULL} $=1000$ per step closer to her head, whatever the goal
            (\code{goal_progress}).
      \item \code{KILL_ENEMY_QUEEN} $=$ \code{false} turns both off.
\end{itemize}
\textbf{Scenario test} (\code{sim/scen_qkill.py}, round 10): their queen died head on head in all three cases.
\begin{itemize}
      \item adjacent: one step;
      \item a 6-long worker two tiles away: one step after she moved towards it;
      \item a 3-long worker three tiles away: two steps, one of them paid.
\end{itemize}
\textbf{Simulator games} (seed 6, 4 founders a team, pearls 30$\times$ sparser; \code{probe_qkill.py}). In this simulator both teams'
founders start next to each other, so a founder of ours traded itself for their queen on round 1. Both sides then had 3 dragons,
but in both games (sides swapped) the killing team fell behind and was wiped out:
\begin{itemize}
      \item game 1: 12 against 26 dragons by round 60;
      \item game 2: 11 against 21;
      \item control without the rule: 18 against 17.
\end{itemize}
Two games, a simulator with the wrong economy and starts unlike the real maps (where teams begin at opposite ends): to be checked in
the sandbox.


\section{The official maps in the simulator}
\code{sim/engine.py --map NAME --seed S} now plays any of the 25 official maps as the official engine does
(\code{config.cc}, \code{pearls.cc}):
\begin{itemize}
      \item \textbf{Map contents}: kelp, portals, starting dragons in file order (ids 0 and 1 are the queens), rounds from 0, the map's
            unit limit (none of the 25 sets one: 64).
      \item \textbf{Pearl beds}: a bed and its mirror share one countdown. The gaps are drawn with the engine's own generator
            (\code{std::mt19937_64} seeded with the match seed), so with a replay's seed every pearl spawns on the same round as in
            that match. Only a spawn blocked by a dragon on the bed can differ.
      \item \textbf{Checks}: every map loads, kelp mirrors on all of them, every portal crossing leads back, and the generator matches
            the C++ standard's reference value.
\end{itemize}
\textbf{Two lost games reproduced} (v23 against itself, so where our pearls go rather than how they play; (dragons, total length)):
\begin{center}
      \begin{tabular}{@{}llllll@{}}
            \toprule
            Map, seed                             &            & round 25     & round 50         & round 75         & round 100        \\
            \midrule
            Schooltime, \code{0xa5e0c0ddde697915} & top team   & (5, 11)      & (7, 18)          & (10, 26)         & (22, 52)         \\
                                                  & v23 before & (4, 8)       & (13--15, 34--37) & (23--24, 62)     & (20--34, 54--84) \\
                                                  & v23 now    & (5, 11)      & (16--18, 41--43) & (24--30, 62--74) & (33--35, 85--87) \\
            \midrule
            Autarky, \code{0xfd11114afed14603}    & top team   & (12, 30)     & (8, 19)          & (16, 38)         & (15, 38)         \\
                                                  & v23 now    & (12, 33--36) & (12--14, 29--34) & (9--12, 28--32)  & (14--15, 38--41) \\
            \bottomrule
      \end{tabular}
\end{center}
Our pearl economy keeps up with theirs on both. What lost Schooltime was the queen (next section).


\section{Boxed in: keep the part that can live}\label{sec:keep}
\textbf{What happened.}
\begin{itemize}
      \item \textbf{Schooltime}: each queen starts 4 long, sealed in a 2$\times$2 kelp box, so every move hits her own body. Ours died
            on round 0. The top team splits off the 2-long tail; that child is boxed too and dies, and the queen eats the pearl it
            leaves. She then lives 3 long, circling the box all game. Their (5, 11) at round 25 is exactly that. Since the queen's
            length is the first tiebreak, the match was lost on round 0.
      \item \textbf{Slithery Fight}: the queen starts coiled with no free side. The split escape always kept 2 (the v13 rule: the rear
            $L-2$ leave, the head part dies), and on round 1 the remnant rule made her die on her own neck.
      \item The same rule cost a unit in most split escapes: the 2-long remnant died even where it could have followed the child out.
\end{itemize}
\textbf{Change} (\code{endgame.hpp}, \code{choose_escape}). The child moves later the same round:
\begin{itemize}
      \item its first step frees its tail, our segment \code{keep};
      \item a child that cannot move at all (boxed, and under 4 so it cannot split) dies at once and frees all of its tiles.
\end{itemize}
When such a tile is beside our head, we live on as well. The split keeps the most length that does this:
\begin{itemize}
      \item Candidate keeps: each body segment beside the head, $L-2$, $L-3$, and 2 (the old escape).
      \item A plan wins on length saved: ours if we live, plus the child's if it has a way out.
      \item The queen (the part we keep) must live first.
      \item A child that dies at once is a price paid only when our head is otherwise certainly dead.
      \item All earlier gates (\code{fine}, room, tail threatened) still come first.
\end{itemize}
Supporting changes:
\begin{itemize}
      \item \textbf{The child holds} (\code{main.cpp}, \code{escape_child_holding}). A child that got its parent's body handover makes
            no ordinary split for its first \code{ESCAPE_CHILD_HOLD_TURNS} $=3$ turns (a split would leave its tail where it is). On its
            first turn, a move that leaves its tail in place (eating with every step) costs \code{ESCAPE_CHILD_STALL_COST} $=5000$.
            The handover now carries the child's segments whatever we keep.
      \item \textbf{Remnant clear} (\code{play_clear_remnant}): only when boxed in (no free side), and never the queen.
      \item \textbf{Ordinary splits check the parent} (\code{split_parent_ok}). The child moves right after the split, from our tail,
            and may take a free side of our head next to it. We need a free side away from the child, or two beside it. On Slithery
            Fight the queen's founder split on round 1 had just one, and the child took it.
      \item \textbf{Our queen's last ways out} (\code{queen.hpp}, \code{queen_exit_block}). She moves first each round. A worker's
            move onto a free side of her head, while she has at most two, costs \code{QUEEN_EXIT_BLOCK_COST} $=4000$.
\end{itemize}
\textbf{Checked} (round-0 to round-2 traces):
\begin{itemize}
      \item \textbf{Schooltime}: split 2, the child dies, she eats its pearl and circles the box at length 3. Both queens are alive at
            rounds 60 and 200.
      \item \textbf{Slithery Fight}: she keeps 7 (the child of 18 leaves first), follows it out, and seeds as usual. Before, she died on
            round 1.
      \item \textbf{Our deaths in rounds 0--3} (team A), before and now:
            \begin{itemize}
                  \item Help 1 and 0;
                  \item Unsw 3 and 2;
                  \item Autarky 3 and 2;
                  \item Slithery Fight 8 and 4.
            \end{itemize}
            The escapes that still lose their remnant have no body segment beside the head, e.g.\ a straight body nosed into a dead end.
      \item \textbf{Small}: still lost on round 3. She starts in a 1-wide corridor that runs into the corner where two of our starting
            dragons are, and that worker had only two moves, both next to her.
\end{itemize}


\section{Farming good beds}
A bed's respawn rate is learnt from the countdowns seen on it (\code{bed_gap}). The countdown runs $g, g-1, \dots, 1$ after each
spawn, so its mean over time is about $(g+1)/2$. With \code{GOOD_BED_MIN_SAMPLES} $=3$ rounds of it seen, a bed whose estimated
mean gap is at most \code{GOOD_BED_GAP} $=12$ is \emph{good}. A good bed is:
\begin{itemize}
      \item an objective \code{GOOD_BED_SOON_ROUNDS} $=8$ rounds before it spawns (other beds: 4);
      \item once its timer has run out while nobody watched, an objective at cost \code{GOOD_BED_EXPIRED_COST} $=2$ for up to
            \code{GOOD_BED_EXPIRED_MAX_AGE} $=150$ rounds (other beds: 5, for 40);
      \item worth 80 instead of 50 in \code{region_value};
      \item free of \code{VISIT_COST} when walked over again (farming is not wandering in circles).
\end{itemize}
\textbf{Effect} (official maps, 100 rounds, new build against the build without it, sides swapped). Bed pearls eaten, both games:
\begin{itemize}
      \item Devil 100 against 76; Schooltime 231 against 198; Portals 121 against 110;
      \item Autarky and Default even; Weakhold not changed at all.
      \item Overall $+12\%$.
\end{itemize}
The limit is memory: every split starts a new process that knows nothing, and most dragons are young.

\textbf{Where the rest goes.} Measured yield is far below what the beds give:
\begin{itemize}
      \item Devil's 18 one-round beds give 1,800 pearls per 100 rounds; we eat about 50.
      \item They lie in 1-wide vertical tubes, and Weakhold's and Stripes' good beds sit in dead ends.
      \item On Portals, most 2$\times$2 bed blocks cannot be reached at all, and the rest are 1-tile dead ends behind portals.
\end{itemize}
Farming those needs a dragon that stays and loops through a tube, a bigger change than this one.

\section{Pocket discipline: not changed}
Raids already need 3 pearls inside the pocket (\code{POCKET_MIN_PEARLS}), and the keep-the-most escape above now saves the unit
when a pocket is 2 wide. In a 1-wide dead end the head part still dies, as before. No new rule: the risk of dragons refusing pockets
outweighs what is left.

\section{v23 against the morning's v23 (simulator, official maps, 200 rounds, sides swapped)}
\begin{center}
      \begin{tabular}{@{}lll@{}}
            \toprule
            Map            & New as A                 & New as B                \\
            \midrule
            Schooltime     & won (queen 3 against 0)  & won (queen 3 against 0) \\
            Slithery Fight & won                      & won                     \\
            Weakhold       & won (queen 14 against 0) & won (queen 6 against 0) \\
            Help           & lost                     & won                     \\
            Devil          & won (wiped them out)     & lost (wiped out)        \\
            Autarky        & won (queen alive)        & lost                    \\
            \bottomrule
      \end{tabular}
\end{center}
8 of 12 by the official ranking. Devil is a wipe-out race whichever side starts first. No crashes or timeouts.


\section{The queen stays out of fights}\label{sec:queensafe}
\textbf{Why.} In self-play 27 of 50 queens were dead by round 100, almost all 2--3 long while seeding. In real games she also:
\begin{itemize}
      \item crossed a portal whose far side held a dragon, and died at once;
      \item was boxed in by our own dragons.
\end{itemize}
\textbf{Changes.}
\begin{itemize}
      \item \textbf{Valuable from round 0} (\code{compute_roles}). Before, she was valuable only from round 75, when she starts
            keeping her length. Now she has the king's caution throughout:
            \begin{itemize}
                  \item enemy reach threats rule moves out;
                  \item crowd and unseen-head risks are priced;
                  \item a portal whose exit we cannot see now is never a safe move;
                  \item no hunting and no head trades (\code{eligible_target});
                  \item escape sprints are allowed.
            \end{itemize}
      \item \textbf{Her death is priced high} (\code{death_cost}): at least \code{QUEEN_DEATH_LENGTH} $=30$ segments, whatever her length.
            Every risk is priced against it. A blind portal now costs her about 9,600 points; at length 3 it was about 1,500.
      \item \textbf{No fight goal} (\code{prepare_fight}) and \textbf{no centre rush} (\code{centre_pull}): the middle is where the
            fights are.
      \item \textbf{Away from enemy heads}: \code{QUEEN_ENEMY_DISTANCE_STEP} $=150$ per step from the nearest one (workers: 12).
      \item \textbf{Room from our own}: a worker's move onto any free side of her head costs \code{QUEEN_SIDE_COST} $=600$, and 4,000
            when she has at most two free sides (\code{queen_exit_block}).
      \item \textbf{She keeps 4 while seeding} (\code{QUEEN_SEED_KEEP}). At 2--3 long she could neither escape-split nor outlast a
            crowd. She now splits only from 6 long, giving the child everything above 4.
\end{itemize}
\textbf{Queen deaths by round 100} (self-play, all 25 maps, seed 1, 50 queens):
\begin{center}
      \begin{tabular}{@{}ll@{}}
            \toprule
            Build                                                & Queens dead \\
            \midrule
            before (this morning's split fixes and farming only) & 27          \\
            cautious, still splitting down to 2                  & 18          \\
            cautious, and nobody hunts queens (control)          & 15          \\
            cautious, keeps 4 while seeding (this build)         & 13          \\
            \bottomrule
      \end{tabular}
\end{center}
Still lost:
\begin{itemize}
      \item \textbf{Small}: both queens on round 3, as before (her corridor runs into a corner our starting dragons fill).
      \item \textbf{Prisoners Dilemma}: both, boxed at length 3. She starts 3 long, so she cannot split there.
      \item The rest are scattered head-on deaths.
\end{itemize}
The portal and own-team cases are covered by these rules; no scenario was rebuilt for them.

\textbf{Against the previous build} (8 official maps, 200 rounds, sides swapped): won 11 of 16. Her queen was alive at the end in 11
games; the old one's in 4.
\begin{itemize}
      \item \textbf{Seeding is unaffected}: 301 against 305 dragons at round 50, summed over the 16 games.
      \item \textbf{Devil} went against us twice. Four more Devil games split 2--2, so that looks like a wipe-out race with no clear
            favourite rather than the change.
      \item Some wins were by the queen alone, with a smaller total (Default 8 against 67, Weakhold 24 against 76).
\end{itemize}



%=====================================================================
\part{v23 $\rightarrow$ v24: portals}
%=====================================================================
v24 starts here. Everything before this part is v23, which ended with the queen staying out of fights
(section~\ref{sec:queensafe}).

\section{Portals: rooms, pockets, and pacing in front of them}\label{sec:portals}
\textbf{Measured} (\code{sim/portal_loiter.py}; v23 against itself, official maps, seed 3, 150 rounds). Two separate faults:
\begin{itemize}
      \item \textbf{Shut in a portal room.} Default's centre (a grid of 4$\times$4 rooms), Autarky's 3$\times$3 boxes and Portals' cells are
            sealed by kelp; the only ways out are portals on their inner edges. Leaving means a crossing whose far tile we cannot
            see, priced at 30\% death (\code{P_BLIND_PORTAL}). A worker that had eaten what was there circled the room instead:
            on Autarky four dragons for 108--142 rounds each, on Default one for 124.
      \item \textbf{Pacing in front of a portal.} The searches behind the goals (pearl route, objective map) treat a crossing as an
            ordinary step, but the move prices it as a 15--30\% death. So the goal led a worker up to the portal, the crossing lost
            (e.g.\ $-1460$ against $-562$ for turning away), the goal led it back. Rounds a dragon spent pacing (at most 4 tiles in
            12 rounds) near portals: 125 on Portals, 24 on Islands. Away from portals: 9.
      \item Most such pulls on Portals came from \textbf{pearls in portal pockets}: 1-tile cells whose only side out is a portal.
            A small dragon that goes in eats and dies there: its neck is in the portal and it is too short to split out.
\end{itemize}
How risky blind crossings really are (\code{sim/portal_exit_risk.py}, share of rounds a crossing's far tile held a dragon):
\begin{itemize}
      \item rounds 0--49: 2--14\%;
      \item rounds 50--149: 8--19\%;
      \item later: 12--25\%.
\end{itemize}
\code{P_BLIND_PORTAL} $=0.3$ is pessimistic but not far off late in a game. It is left as it is for the fixes in this section;
section~\ref{sec:portaltune} then prices crossings by round.

\textbf{Changes} (\code{main.cpp}, \code{scoring.hpp}):
\begin{itemize}
      \item \textbf{Shut in a room} (\code{update_portal_room}, once a turn). We are shut in when:
            \begin{itemize}
                  \item the tiles we can reach without a portal (at most \code{PORTAL_ROOM_MAX_CELLS} $=64$) are all seen;
                  \item they hold no objective at all;
                  \item a portal leads out.
            \end{itemize}
            Then a crossing (not onto a body we know of) is priced at \code{P_ROOM_EXIT} $=0.1$ and is worth \code{ROOM_EXIT_VALUE}
            $=1500$. A valuable dragon may then cross blind too (its safety filter allows it), and a portal we never looked through counts
            as exploring. Not the queen: a room is a fine place for her.
      \item \textbf{Portal pockets are not objectives} (\code{portal_pocket_worth}). A pearl or bed whose portal-free tiles are few
            (at most \code{PORTAL_POCKET_MAX_CELLS} $=8$), all known, and a dead end or too small to turn in, is left off the
            objective map. The exception is when a pocket raid is allowed and it holds \code{POCKET_MIN_PEARLS}, as for kelp pockets.
      \item \textbf{A portal on the way is crossed or dropped} (\code{refuse_route_portals}). The crossing's candidate records whether
            it makes progress on the goal (\code{route_portal}). If we then choose another first step, its far tile is marked a trap
            for \code{PORTAL_REFUSE_ROUNDS} $=15$, so the route and objective searches go round it and the goal picks something else.
      \item The objective map now records which tiles are objectives themselves (\code{objective_src_round_}).
\end{itemize}
\textbf{Results} (same probe):
\begin{center}
      \begin{tabular}{@{}lll@{}}
            \toprule
                                                       & Before               & After         \\
            \midrule
            Autarky: stays $\geq$20 rounds by a portal & 7 (four of 108--142) & 2 (26 and 22) \\
            Default: stays $\geq$20 rounds             & 15                   & 1             \\
            Portals: rounds pacing near portals        & 125                  & 5             \\
            Islands: rounds pacing near portals        & 24                   & 1             \\
            \bottomrule
      \end{tabular}
\end{center}
\textbf{Against v23} (8 portal maps, 200 rounds, sides swapped): 9--7.
\begin{itemize}
      \item Maze, Prisoners Dilemma and Stripes were decided by side, not build: mirrored results, or the A side won every game
            (Stripes, 4 more seeds: 4--4).
      \item Total length 1,198 against 1,362, most of it one Schooltime game (93 against 393).
\end{itemize}
\textbf{Points} (metered wasm, 12,345 recorded turns on Unsw, 64$\times$64): median 17.3M $\rightarrow$ 17.5M, max 43.3M
$\rightarrow$ 43.6M.


\section{Open early, only good ones late}\label{sec:portaltune}
Many portal maps keep pearls on the far side of their portals, and early in a game whoever gets through first eats them. The
crossing risk was 30\% at any time; the simulator measured 2--14\% in rounds 0--49 and 12--25\% after round 150
(section~\ref{sec:portals}). So:
\begin{itemize}
      \item \textbf{A blind crossing's risk follows the round} (\code{blind_portal_p}): \code{P_BLIND_EARLY} $=0.08$ up to round
            \code{PORTAL_RISK_EARLY_ROUND} $=40$, rising linearly to \code{P_BLIND_LATE} $=0.20$ at \code{PORTAL_RISK_LATE_ROUND}
            $=150$. Half that when we saw the far tile free in the last 4 rounds (was a flat 0.1). The constants
            \code{P_BLIND_PORTAL} and \code{P_SEEN_FREE_PORTAL} are gone.
      \item \textbf{The early window is longer}: \code{PORTAL_EXPLORE_END_ROUND} $80 \rightarrow 100$. Until then, as before, workers
            treat a portal they have never seen through as an exploration objective and an ordinary move (5\% risk).
      \item \textbf{Late, only good portals.} From round 100 the far side is valued by the pearls and timers we know of there
            (\code{region_value} with unexplored tiles at 0): \code{PORTAL_GOOD_PER_PEARL} $=300$ per expected pearl more than
            around us, up to \code{PORTAL_GOOD_CAP} $=2500$.
            \begin{itemize}
                  \item Nothing for a far tile with no way on (a pocket, a dead end), unless it is a planned pocket raid.
                  \item Before round 100 the far side keeps its old small weight (48 per pearl, unexplored tiles counted, $\pm$360).
                  \item For a 3-long worker at round 150 the crossing's risk is about 1,000, so it needs about 3 more known pearls on
                        the far side; a longer dragon needs more.
            \end{itemize}
\end{itemize}
\textbf{Tried and dropped} (each against the build with only the portal fixes above, 200 rounds, sides swapped):
\begin{itemize}
      \item A bonus of 720 for an early crossing the goal's way goes through lost all 4 Default games (total length 68 against 238).
      \item Valuing the far side at 500 per pearl from round 0, unexplored tiles counted, went 4--8 and lost all 4 Default games.
            On Default's grid of portal rooms both sent the swarm hopping room to room in the middle of the map.
\end{itemize}
\textbf{Results.}
\begin{itemize}
      \item \textbf{Against the portal fixes alone} (6 maps, sides swapped): 7--5, total length 1,479 against 1,099.
            Help, Islands and Portals went to the B side in both orders.
      \item \textbf{Early crossings} (rounds 0--99, seed 5, the team that crossed more): Portals $12 \rightarrow 53$, Default $28 \rightarrow 78$,
            Islands $58 \rightarrow 111$.
      \item \textbf{The cost}: more small dragons die finding out where an unknown portal leads. On Portals many lead into
            1-tile pockets, where a 2--3-long dragon eats and dies. Deaths on a crossing move, early: Portals $4 \rightarrow 11$,
            Islands $7 \rightarrow 23$.
\end{itemize}
\textbf{v24 against v23} (all of this part against the build that ended v23; 6 portal maps, 200 rounds, sides swapped): 5--7,
total length 1,275 against 1,201. Seven of the 12 were settled by whose queen was alive or longer, not by length: Autarky as A
had 122 against 17 and lost on the queen (4 against 11). Read it as even. Since the queen tiebreak swamps everything else in
these runs, the portal changes are best judged by the loitering and crossing figures above, and by sandbox games on portal maps.

%=====================================================================
\part{v24 $\rightarrow$ v25: dispersion}
%=====================================================================
v25 starts here. Everything before this part is v24.

\section{Spread out early}\label{sec:disperse}
\textbf{The problem.} We keep losing early eliminations, and losing consistently means exploration is not luck. Several of our
workers going to the same spot gain little, since one or two there split quickly and fill it. They should spread over the map
instead.

\textbf{What was there already:}
\begin{itemize}
      \item \textbf{Pearls}: each pearl is left to whichever known teammate reaches it first (\code{pearl_taken_by_ally}: visible
            teammates by path steps, others by the position they last sent, at most a round old).
      \item \textbf{Destinations} (v14, \code{explore.hpp}): every worker heads for a part of the map nobody has explored, chosen
            by its id. Only on maps of 1,600 tiles or more and from round 20, so not on most official maps.
      \item \textbf{Exploring}: nothing shared it out. Every never-seen tile within 30 steps was an objective for everyone, so dragons
            near each other headed for the same unexplored ground, and on the centre-rush maps every worker ran to the middle.
\end{itemize}
\textbf{Measured} (\code{sim/spread.py}, v24 against itself): teammate heads within 6 tiles of each head, against what an even
spread over the map would give.
\begin{itemize}
      \item Trophy, round 50: 9.6 against 4.3.
      \item Schooltime, round 75: 3.7 against 1.6.
      \item Default, round 50: 3.3 against 1.8.
\end{itemize}

\textbf{Changes} (\code{main.cpp}, \code{scoring.hpp}):
\begin{itemize}
      \item \textbf{Shared exploring} (\code{ally_claims}, in \code{build_objective_map}). Until \code{CLAIM_END_ROUND} $=200$, a
            never-seen tile, or a timer that ran out while nobody watched, that a known teammate reaches first costs
            \code{CLAIMED_EXTRA} $=10$ more (a never-seen tile costs 8).
            \begin{itemize}
                  \item Known teammates: those we see, by their path steps (\code{ally_etas_}), and others by the position they sent
                        in the last \code{CLAIM_SIGNAL_AGE} $=4$ rounds plus that age.
                  \item Ties go to the lower id.
                  \item It is a cost, not a ban: with nothing else to do, a dragon still explores ground a teammate is nearer to.
            \end{itemize}
      \item \textbf{At most two rush the centre} (\code{centre_crowded}). On the centre-rush maps a dragon with
            \code{CENTRE_MAX_RUSHERS} $=2$ or more known teammates closer to the centre gets no centre pull.
\end{itemize}

\textbf{Tried and dropped} (each against v24, the same 6 maps, 150 rounds, sides swapped):
\begin{itemize}
      \item \textbf{Destinations on every map from round 0} (the v14 rule without its size and round limits): 3--9, and 24\% less
            length at round 50 (22.8 against 30.0). Sending workers away from food they can see costs more than it explores, as v14
            found.
      \item \textbf{Stronger shared exploring} (extra cost 20, positions up to 6 rounds old): 6--6, against 8--4 for the settings
            kept.
\end{itemize}

\textbf{Results against v24} (12 official maps, 24 games, 150 rounds, sides swapped; \code{sim/ab.sh}, \code{sim/ab_econ.py}):
\begin{itemize}
      \item \textbf{Games}: 14--10 by the official ranking.
      \item \textbf{Wipe-outs}: we wiped them out 6 times; they wiped us out once.
      \item \textbf{At round 50}, on average: 9.7 dragons and 24.5 length, against 8.4 and 22.2.
      \item Several of the 10 losses were on the queen tiebreak with more length (e.g.\ Portals as A, 47 length against 39 at round 100).
\end{itemize}
\textbf{Points} (metered wasm, 11,952 recorded turns on Unsw, 64$\times$64): median 17.4M $\rightarrow$ 18.5M, max 46.7M
$\rightarrow$ 46.8M.


\section{Anticipated pearls}\label{sec:forecast}
Every tile in view shows how many rounds until its next pearl, and the bot already remembered that countdown
(\code{CellMemory}). Until now it only used it within 4 rounds of the spawn (10 for the king), at cost rounds left $+1$ plus the
walk. So a dragon passing a bed knew exactly when it would fill, and forgot about it until it was almost due.

\textbf{Timed beds on the official maps} (beds that cannot spawn before round 20 or later):
\begin{itemize}
      \item \textbf{Prisoners Dilemma}: the two middle columns (32 beds) all spawn once, around rounds 479--489. 44 beds either side of
            the middle first spawn around 150--250, then every 150--250 rounds.
      \item \textbf{Tower Defense}: 38 beds (centre and corners) spawn once, between rounds about 399 and 499.
      \item \textbf{Unsw}, \textbf{Maze}: 36 and 28 beds at 250--500. \textbf{Australia}: 126 beds starting after round 40 or 100.
            \textbf{Schooltime}: nothing before round 19.
      \item \textbf{Queen of Spades But She Ages}: its 170 delayed beds (rows timed 40--80 at the edges up to 400--480 near the middle)
            are all sealed off by kelp (\code{sim/engine.py} reachability from every start). They spawn but nobody can eat them,
            exactly as on the normal Queen of Spades, where the same tiles are 1--1000 beds.
\end{itemize}

\textbf{Change} (\code{main.cpp}, \code{spawn_forecast_cost}). Any map, no map table:
\begin{itemize}
      \item \textbf{When to set off.} A known future spawn becomes an objective once its rounds left are at most our turns to get
            there plus \code{FORECAST_MAX_WAIT} $=2$, so we arrive as it spawns.
      \item \textbf{Its cost} is that wait, max(0, rounds left $+\,1\,-$ our turns there); the objective search adds the
            walk. A pearl that is about to appear right where we can be in time ranks like one already there.
      \item \textbf{Bursts rank higher.} Each other spawn within \code{FORECAST_CLUSTER_RADIUS} $=3$ tiles and
            \code{FORECAST_CLUSTER_WINDOW} $=6$ rounds takes $1/$\code{FORECAST_CLUSTER_DIV} $=1/3$ round off the cost, at most
            \code{FORECAST_CLUSTER_CAP} $=8$.
      \item The objective radius bounds it: 30 steps for workers, 200 for the king.
      \item \code{SPAWN_SOON_ROUNDS}, \code{KEEPER_SPAWN_SOON_ROUNDS} and \code{GOOD_BED_SOON_ROUNDS} are gone.
\end{itemize}
\textbf{Tried and dropped}: making every known future spawn an objective at once (cost = the wait). Dragons went early and circled
on the beds; a dragon's body on a bed when it spawns blocks the pearl. In a test copy of Prisoners Dilemma with its burst moved to
round 80 (\code{sim/timed_beds.py}, 3 seeds, sides swapped) that build ate 59 of the burst's pearls against 79 for v25 without
forecasting. Setting off on time turned it round: 77 against 60.

\textbf{Results} (against v25 without forecasting):
\begin{itemize}
      \item \textbf{12 official maps}, 150 rounds, sides swapped: 14--10. At round 50, 9.0 dragons and 23.8 length against 8.5 and
            23.0.
      \item \textbf{Prisoners Dilemma}, 4 seeds, full games, sides swapped: won 7 of 8. The delayed beds' pearls eaten: 126 against
            59. In self-play these games end by wipe-out between rounds 196 and 277, so the round-480 burst itself never came.
      \item \textbf{Tower Defense}, 2 seeds: 2--2, all decided by the A side wiping the other out before the round-400 beds.
      \item \textbf{Points} (metered wasm, 23,150 recorded turns on Big Empty, 64$\times$64, every tile a bed): median 18.3M
            $\rightarrow$ 18.4M, max unchanged at 29.1M.
\end{itemize}

%=====================================================================
\part{v25.1.0: the queen waits by splitting}
%=====================================================================
\textbf{Implemented: 6 October 2026.} Base: \code{bots/New/current/bot-v25.0.0-omar}.
The new version is \code{bots/New/current/bot-v25.1.0}; the base folder is unchanged.

\section{Keep the queen's length while blocked}\label{sec:queenwait251}
\textbf{The problem.} A queen whose head is blocked may survive if the blocking dragon moves away on a later turn.
The existing escape could hand off most of her body to a child and leave only a two-segment queen, giving her no further
split with which to wait. v25.1.0 instead spends the smallest legal amount of length to buy another turn.

\textbf{Change} (\code{endgame.hpp}, \code{play_split_escape}):
\begin{itemize}
      \item The existing queen emergency test is unchanged: there are no candidate moves, or all candidates are judged hopeless
            by the existing head-trade, exit and survival-search checks. Ordinary movement is still preferred when this test finds
            a way on; this is not a new general-purpose waiting rule.
      \item When the queen has at least four segments and the existing split-slot checks allow it, she issues \code{SPLIT 2}.
            The two tail segments become a child; the queen keeps her original head position, identity and remaining $L-2$ segments.
            Splitting consumes her action for that turn.
      \item On the next turn, normal movement assessment runs again. She repeats the two-segment split while the emergency
            condition remains true and a split remains available, and resumes movement when a viable move becomes available.
      \item The wait split does not require a complete body trace or an escape route for the child. It can therefore buy the queen
            a turn even if the child is boxed and dies. The queen is not marked as a sacrificial escape remnant; the existing
            status and body-handover messages are still sent.
\end{itemize}

\textbf{Limits.} Two segments is the smallest legal child; the originally suggested one-segment child is not legal.
The queen must retain at least two segments, so this behavior stops at length two or three. Existing team-cap and
escape-slot reservations are unchanged. Waiting can help a temporary blockage, but does not guarantee escape from a permanent box.

\textbf{Scope.} Only \code{endgame.hpp} changes executable behavior. \code{queen.hpp} updates the description of the queen's
escape. Normal queen seeding, movement and cycling, worker splits and escapes, scoring, routing and all configuration settings
are unchanged from Omar's base.

\section{Verification of v25.1.0}
\textbf{Focused tests: 53 checks passed.} \code{battlecode-automation/scripts/test_v251.py} runs the real metered wasm bots
and the actual game engine using \code{unswbc 1.2.9}. Coverage includes both queen IDs, all queen phases, lengths 4/5/8/10,
repeated splits, complete and incomplete body knowledge, a boxed child, release of a blocking ally, minimum-size and unit-cap
boundaries, reserved slots, and matching baseline actions for ordinary queen and worker behavior.

\textbf{Observed engine example} (round numbers start at zero):
\begin{center}
      \begin{tabular}{@{}rllr@{}}
            \toprule
            Round & Queen head before action & Action & Length before action \\
            \midrule
            0 & $(8,8)$ & \code{SPLIT 2} & 8 \\
            1 & $(8,8)$ & \code{SPLIT 2} & 6 \\
            2 & $(8,8)$ & \code{MOVE E}  & 4 \\
            3 & $(9,8)$ & Movement continues & 4 \\
            \bottomrule
      \end{tabular}
\end{center}
The ally clears the exit after the queen's turn in round 1. She survives and moves on round 2. In the same straight-body
fixture, Omar's base issues \code{SPLIT 6}, retains only two segments and dies before the blocker clears. A separate fixture
keeps the head blocked longer and verifies three consecutive waits: $8\rightarrow6\rightarrow4\rightarrow2$.

\textbf{Full-match sandbox check against Omar's base:} six completed 500-round games, on Default Small, Schooltime and
Slithery Fight, each from both sides with seed 25101. Results: \textbf{3 wins, 0 draws, 3 losses}; zero evaluation errors and
zero bot runtime errors across both versions. The candidate's highest measured turn cost was 94,449,969 of the 100,000,000-point
budget. This small sample does not establish a win-rate improvement or a reduction in elimination losses; the focused engine
fixtures establish that the requested waiting behavior works.

\textbf{Evidence.} The test output, source/map hashes, match ledger and detailed validation are under:
\par{\small\raggedright\path{battlecode-automation/experiments/v251/queen-wait-smoke/}\par}
The directory includes \code{behavior-tests.txt}, \code{manifest.json}, \code{matches.jsonl} and \code{VALIDATION.md}.
No new strategy tuning was included in this version.

\end{document}
